\documentclass[11pt,a4paper,twoside,openright]{book}
% Shared preamble for the Derivatives training sheet.
\usepackage[T1]{fontenc}
\usepackage[utf8]{inputenc}
\usepackage{lmodern}
\ifdefined\KindleEdition
  % 6-inch e-ink page (Kindle Paperwhite): the page *is* the screen, so the
  % margins are only what the eye needs, not what a printer needs.
  \usepackage[paperwidth=90mm,paperheight=122mm,
              top=4mm,bottom=6mm,left=4.5mm,right=4.5mm]{geometry}
\else
  \usepackage[a4paper,margin=2.4cm]{geometry}
\fi
\usepackage{amsmath,amssymb,amsthm}
\usepackage{microtype}
\usepackage{enumitem}
\usepackage{xcolor}
\usepackage[most]{tcolorbox}
\usepackage{fancyhdr}
\usepackage[hidelinks]{hyperref}

\definecolor{accent}{HTML}{1F4E79}
\definecolor{soft}{HTML}{F4F7FA}

% --- function graphs ------------------------------------------------------
\usepackage{pgfplots}
\pgfplotsset{compat=1.18}
\usetikzlibrary{arrows.meta}

% One consistent look for every function graph in these books.
\pgfplotsset{
  fnaxis/.style={
    width=5.9cm, height=4.6cm,
    axis lines=middle,
    axis line style={-{Stealth[length=4pt]}, gray!65!black, line width=0.3pt},
    tick style={gray!65!black},
    tick label style={font=\tiny, gray!50!black},
    label style={font=\scriptsize},
    xlabel={$x$}, ylabel={$y$},
    every axis x label/.style={at={(ticklabel* cs:1.0)}, anchor=west},
    every axis y label/.style={at={(ticklabel* cs:1.0)}, anchor=south east,
      xshift=-1pt},
    every axis plot/.append style={accent, very thick, smooth, samples=201},
    clip=true, enlargelimits=false,
    title style={font=\scriptsize, yshift=4pt},
  },
  fnaxis/.append style={/pgfplots/legend style={font=\tiny, draw=none,
    fill=none, at={(0.98,0.03)}, anchor=south east}},
}

% A row of small graphs, evenly spaced with \hfill between them, plus an
% optional caption line: \begin{gallery} <pic>\hfill<pic> \gcap{...} \end{gallery}
\newenvironment{gallery}
  {\par\medskip\noindent\ignorespaces}
  {\par\medskip\noindent}
\newcommand{\gcap}[1]{\par\smallskip{\centering\footnotesize #1\par}}

% --- plane regions, implicit curves, surfaces -----------------------------
\usepgfplotslibrary{fillbetween}
\definecolor{regionfill}{HTML}{9DB4D8}
\pgfplotsset{colormap={house}{color=(soft) color=(accent!45) color=(accent)}}

\pgfplotsset{
  % Squared, gridded axis for plane regions and implicitly given curves.
  geoaxis/.style={
    fnaxis,
    width=5.9cm, height=5.9cm,
    grid=both,
    grid style={gray!22, line width=0.2pt},
    minor tick num=1,
  },
  % Level-curve maps: many thin curves, no single curve emphasised.
  contouraxis/.style={
    geoaxis,
    every axis plot/.append style={accent!75, line width=0.5pt, smooth,
      samples=201},
  },
  % Surfaces: \addplot3[surf, domain=..., y domain=...] {expr};
  surfaxis/.style={
    width=6.0cm, height=5.0cm,
    view={130}{22},
    colormap name=house,
    ticks=none,
    xlabel={$x$}, ylabel={$y$}, zlabel={$f$},
    label style={font=\tiny, gray!50!black},
    title style={font=\scriptsize},
    every axis plot/.append style={samples=25, shader=faceted interp,
      draw=accent!35, line width=0.1pt},
  },
}

% Shorthand for a shaded region between two named paths.
% \region{<path A>}{<path B>}{<soft clip options>}
\newcommand{\region}[3]{%
  \addplot[regionfill, opacity=0.55, forget plot, sharp plot]
    fill between[of=#1 and #2, #3];}

% --- rule boxes -----------------------------------------------------------
\newtcolorbox{rulebox}[1]{%
  breakable, enhanced, colback=soft, colframe=accent,
  boxrule=0.4pt, left=7pt, right=7pt, top=6pt, bottom=6pt,
  fonttitle=\bfseries, coltitle=white, colbacktitle=accent,
  title={#1}}

% --- worked examples ------------------------------------------------------
% One fully narrated solution per chapter, carrying the chapter's central
% technique and the mistakes that technique invites:
%   \begin{worked}{Task} ... \wstep{Step 1.} ... \wtrap{...} ... \end{worked}
\newtcolorbox{worked}[1]{%
  breakable, enhanced, colback=white, colframe=accent!65,
  boxrule=0.4pt, left=7pt, right=7pt, top=6pt, bottom=6pt,
  fonttitle=\bfseries, coltitle=accent, colbacktitle=soft,
  title={#1}}

% Run-in step heading inside a worked example.
\newcommand{\wstep}[1]{\par\smallskip\noindent\textbf{#1}\enspace\ignorespaces}

% A pitfall, set off from the running solution so it can be found again.
\newcommand{\wtrap}[1]{%
  \par\smallskip
  {\leftskip=1.3em\noindent
   \llap{\textcolor{accent}{\rule[0.22ex]{0.9em}{0.6pt}}\hspace{0.4em}}%
   \textbf{Trap.}\enspace\itshape #1\par}%
  \smallskip}

% --- problem / solution lists --------------------------------------------
\newlist{problems}{enumerate}{1}
\setlist[problems]{label=\textbf{\arabic*.}, leftmargin=2.1em,
  itemsep=0.55ex, parsep=0pt, topsep=0.8ex}
\newlist{solutions}{enumerate}{1}
\setlist[solutions]{label=\textbf{\arabic*.}, leftmargin=2.1em,
  itemsep=1.0ex, parsep=0pt, topsep=0.8ex}

% --- running heads --------------------------------------------------------
\pagestyle{fancy}
\fancyhf{}
\fancyhead[LE,RO]{\thepage}
\fancyhead[RE]{\small\itshape\nouppercase{\leftmark}}
\fancyhead[LO]{\small\itshape\nouppercase{\rightmark}}
\setlength{\headheight}{14pt}
\renewcommand{\headrulewidth}{0.3pt}
\renewcommand{\chaptermark}[1]{\markboth{\chaptername\ \thechapter.\ #1}{}}
\renewcommand{\sectionmark}[1]{\markright{\thesection\ #1}}

% --- shorthands used throughout -------------------------------------------
\newcommand{\dd}[2]{\frac{\mathrm{d}#1}{\mathrm{d}#2}}
\newcommand{\pd}[2]{\frac{\partial #1}{\partial #2}}
\newcommand{\pdd}[3]{\frac{\partial^{2} #1}{\partial #2\,\partial #3}}
\newcommand{\R}{\mathbb{R}}
\newcommand{\grad}{\nabla}
\DeclareMathOperator{\sgn}{sgn}
\DeclareMathOperator{\Hess}{H}

% German gloss, used on first mention of a technical term
\newcommand{\de}[1]{\emph{#1}}

% --- cross-references to chapters ------------------------------------------
% Chapter files are shared by the standalone volumes and the compendium, so a
% hard-coded "Chapter 6" is wrong in one of them. Write \chref{derivatives}{6}
% instead: it prints "Chapter 6" inside the same volume, "Volume 5, Chapter 6"
% from another volume, and the running number ("Chapter 34") in the
% compendium. \chrefs{vol}{n} is the abbreviated form (Ch. / Vol., Ch.),
% \chnum{vol}{n} the bare number in the current build, and \volref{vol} names
% the volume ("Volume 5", or "Part V" in the compendium). Each volume's
% main.tex sets \thisvolume; the compendium's sets \incompendiumtrue.
\newif\ifincompendium
\newcommand*{\thisvolume}{}
\makeatletter
\@namedef{chref@vol@algebra}{1}\@namedef{chref@off@algebra}{0}
\@namedef{chref@vol@functions}{2}\@namedef{chref@off@functions}{9}
\@namedef{chref@vol@matrices}{3}\@namedef{chref@off@matrices}{14}
\@namedef{chref@vol@foundations}{4}\@namedef{chref@off@foundations}{24}
\@namedef{chref@vol@derivatives}{5}\@namedef{chref@off@derivatives}{28}
\@namedef{chref@vol@integrals}{6}\@namedef{chref@off@integrals}{42}
\@namedef{chref@vol@optimization}{7}\@namedef{chref@off@optimization}{50}
\@namedef{chref@vol@linprog}{8}\@namedef{chref@off@linprog}{60}
\@namedef{chref@vol@mixed}{9}\@namedef{chref@off@mixed}{68}
\newcommand*{\chref@check}[1]{%
  \@ifundefined{chref@vol@#1}{\PackageError{chref}{Unknown volume `#1'}{}}{}}
\newcommand*{\chnum}[2]{\chref@check{#1}%
  \ifincompendium\the\numexpr\csname chref@off@#1\endcsname+#2\relax
  \else#2\fi}
% \chref@same{vol}{same volume}{other volume}
\newcommand*{\chref@same}[3]{\def\chref@tmp{#1}%
  \ifx\chref@tmp\thisvolume\expandafter\@firstoftwo
  \else\expandafter\@secondoftwo\fi{#2}{#3}}
\newcommand*{\chref}[2]{\chref@check{#1}%
  \ifincompendium Chapter~\chnum{#1}{#2}%
  \else\chref@same{#1}{Chapter~#2}%
    {Volume~\csname chref@vol@#1\endcsname, Chapter~#2}\fi}
\newcommand*{\chrefs}[2]{\chref@check{#1}%
  \ifincompendium Ch.~\chnum{#1}{#2}%
  \else\chref@same{#1}{Ch.~#2}%
    {Vol.~\csname chref@vol@#1\endcsname, Ch.~#2}\fi}
\newcommand*{\volref}[1]{\chref@check{#1}%
  \ifincompendium Part~\@Roman{\csname chref@vol@#1\endcsname}%
  \else Volume~\csname chref@vol@#1\endcsname\fi}
\makeatother

% --- Kindle edition: compact layout for a 6-inch screen -------------------
% Everything below applies only when \KindleEdition is defined on the
% command line; the A4 edition is untouched.
\ifdefined\KindleEdition
  \usepackage{etoolbox}
  \usepackage{adjustbox}
  \usepackage{titlesec}

  % The 8cm measure leaves few break points; let TeX loosen a tight line
  % rather than push it into the margin.
  \setlength{\emergencystretch}{1.5em}

  % No running heads: on a 90mm page the header costs more than it gives.
  % A centred page number in the foot is kept so the contents list means
  % something.
  \pagestyle{plain}
  \fancypagestyle{plain}{\fancyhf{}\fancyfoot[C]{\small\thepage}%
    \renewcommand{\headrulewidth}{0pt}}

  % Compact headings: the book class leaves a third of a page above a
  % chapter title, which is unaffordable here.
  \titleformat{\chapter}[display]
    {\normalfont\Large\bfseries\color{accent}}{\normalsize\chaptertitlename\ \thechapter}{0.4em}{\Large}
  \titlespacing*{\chapter}{0pt}{0pt}{1.2em}
  \titlespacing*{\section}{0pt}{1.4ex plus .3ex}{0.7ex}
  \titlespacing*{\subsection}{0pt}{1.2ex plus .2ex}{0.5ex}

  % Boxes: trim the padding, which was set for a 16cm measure.
  \tcbset{left=3.5pt,right=3.5pt,top=3pt,bottom=3pt,
          fonttitle=\bfseries\small,
          halign title=flush left}

  % Problem and solution lists: reclaim the wide left margin.
  \setlist[problems]{label=\textbf{\arabic*.}, leftmargin=1.5em,
    itemsep=0.45ex, parsep=0pt, topsep=0.6ex}
  \setlist[solutions]{label=\textbf{\arabic*.}, leftmargin=1.5em,
    itemsep=0.8ex, parsep=0pt, topsep=0.6ex}

  % Figures and wide tables shrink to the measure instead of running off it.
  % adjustbox only scales down: anything already narrow is left alone.
  \BeforeBeginEnvironment{tikzpicture}{\begin{adjustbox}{max width=\linewidth}}
  \AfterEndEnvironment{tikzpicture}{\end{adjustbox}}
  \BeforeBeginEnvironment{tabular}{\begin{adjustbox}{max width=\linewidth}}
  \AfterEndEnvironment{tabular}{\end{adjustbox}}

  % Display maths was written for a 16cm measure. On 8cm the long chained
  % equations would run off the screen, so any display wider than the text
  % block is scaled down to fit; narrower ones are left at full size. The
  % whole corpus uses \[ ... \] and no other display environment, so this
  % one redefinition covers every case.
  \newsavebox{\kndbox}
  \renewcommand{\[}{\par\addvspace{0.7ex}%
    \begin{lrbox}{\kndbox}$\displaystyle}
  \renewcommand{\]}{$\end{lrbox}%
    {\centering
      \ifdim\wd\kndbox>\linewidth
        \resizebox{\linewidth}{!}{\usebox{\kndbox}}%
      \else
        \usebox{\kndbox}%
      \fi
     \par}%
    \addvspace{0.7ex}}

  % A narrow measure makes perfect justification impossible; prefer a loose
  % line to a line that runs into the margin.
  \sloppy
  \raggedbottom
\fi

\renewcommand*{\thisvolume}{foundations}

\title{\Huge Foundations of Analysis\\[0.4em]
  \large A Training Sheet with Exercises and Worked Solutions\\[0.8em]
  \normalsize Sequences \(\cdot\) Series \(\cdot\) Financial Mathematics
  \(\cdot\) Limits \& Continuity}
\author{}
\date{}

\begin{document}
\frontmatter
\maketitle

\chapter*{About this sheet}
\addcontentsline{toc}{chapter}{About this sheet}

This sheet drills the analysis foundations that a first course in
mathematics for economists opens with, and that every later volume silently
assumes: sequences and their limits, series with the geometric series at
the centre, the financial mathematics built on the geometric series
(compound interest, discounting, present value, annuities), and finally
limits and continuity of functions, up to the Intermediate Value Theorem.

The material is deliberately computational: convergence is established with
the standard limit rules rather than $\varepsilon$-arguments, and every
financial formula is derived once from the geometric series and then
drilled numerically.

The structure is the same in every chapter: \textbf{Core Rules} states the
rules compactly; \textbf{Practice Problems} drills them; \textbf{Solutions}
works every problem through in full. Each chapter closes with
\emph{additional problems} for a second pass.

\tableofcontents

\mainmatter
\chapter{Sequences and Their Limits}

A capital of $1000$ euro earning $2\%$ interest per period does not grow in
one smooth sweep; it grows step by step: $1000$, $1020$, $1040.40,\dots$,
one number per period. Economic processes are full of such step-by-step
lists: iterated interest, instalment plans, prices adjusting period by
period towards an equilibrium, GDP figures year by year. The mathematical
object behind all of them is the \emph{sequence}, and the single most
important question about a sequence is what happens \emph{in the long run}:
does it settle down to a limit, or does it run off beyond every bound?

This chapter builds the toolkit. We state once, informally, what
convergence means, and then compute limits the practical way: with the
limit rules and a short list of standard limits that cover almost every
sequence met in economics. The star among them, $(1+x/n)^{n}\to\mathrm{e}^{x}$,
is exactly the step from period-by-period to \emph{continuous} compounding
and will carry the financial mathematics of a later chapter.

\section{Core Rules}

\begin{rulebox}{Sequences: explicit and recursive}
A \textbf{sequence} is a function that assigns to every index
$n\in\mathbb{N}$ a real number $a_{n}$; we write $(a_{n})_{n\geq 1}$ or
simply $a_{1},a_{2},a_{3},\dots$ It can be given
\begin{itemize}[leftmargin=1.6em,itemsep=0.3ex]
  \item \textbf{explicitly}, by a formula in $n$: e.g.\ $a_{n}=\dfrac{1}{n}$
        or $a_{n}=1000\cdot(1.02)^{n}$ (capital after $n$ periods at
        $2\%$ interest);
  \item \textbf{recursively}, by a start value and a rule for the next
        term: e.g.\ $a_{1}=1000$, $a_{n+1}=1.02\,a_{n}+100$ (interest plus
        a deposit of $100$ each period).
\end{itemize}
\end{rulebox}

\begin{rulebox}{Monotone and bounded sequences}
$(a_{n})$ is \textbf{increasing} if
$a_{n+1}\geq a_{n}$ for all $n$, and \textbf{decreasing} if
$a_{n+1}\leq a_{n}$ for all $n$ (with strict inequalities: strictly
increasing/decreasing). Either way the sequence is called
\textbf{monotone}.

$(a_{n})$ is \textbf{bounded above} if $a_{n}\leq M$ for some fixed $M$ and
all $n$, \textbf{bounded below} if $a_{n}\geq m$ for some fixed $m$, and
\textbf{bounded} if both hold.

\emph{Practical test for monotonicity:} compute the difference
$a_{n+1}-a_{n}$ and determine its sign; for sequences with positive terms
one may instead compare the ratio $a_{n+1}/a_{n}$ with $1$.
\end{rulebox}

\begin{rulebox}{Convergence and the limit}
$(a_{n})$ \textbf{converges} to the \textbf{limit} $a$ if the terms come arbitrarily close to $a$ and stay
there: for every tolerance $\varepsilon>0$, all but finitely many terms
satisfy $|a_{n}-a|<\varepsilon$. We write
\[
  \lim_{n\to\infty}a_{n}=a \qquad\text{or}\qquad a_{n}\to a .
\]
A sequence has \emph{at most one} limit. A sequence that does not converge
is \textbf{divergent}; if its terms eventually exceed every bound we say it
\emph{diverges to $\infty$} and write $a_{n}\to\infty$ (analogously
$-\infty$). In this book we never chase $\varepsilon$'s: limits are
computed with the rules and standard limits below.
\end{rulebox}

\begin{rulebox}{Limit rules}
Suppose $a_{n}\to a$ and $b_{n}\to b$. Then
\begin{itemize}[leftmargin=1.6em,itemsep=0.3ex]
  \item \textbf{Sum rule:} $a_{n}+b_{n}\to a+b$ \ (and $c\,a_{n}\to c\,a$
        for any constant $c$);
  \item \textbf{Product rule:} $a_{n}\,b_{n}\to a\,b$;
  \item \textbf{Quotient rule:} $\dfrac{a_{n}}{b_{n}}\to\dfrac{a}{b}$,
        provided $b\neq 0$;
  \item \textbf{Sandwich rule:} if $a_{n}\leq c_{n}\leq b_{n}$ from some
        index on and $a=b$, then $c_{n}\to a$ as well.
\end{itemize}
\emph{Working method for ratios of polynomials in $n$:} divide numerator
and denominator by the highest power of $n$ that occurs, then apply the
rules together with $1/n^{k}\to 0$.
\end{rulebox}

\begin{rulebox}{Standard limits}
\begin{itemize}[leftmargin=1.6em,itemsep=0.3ex]
  \item $\dfrac{1}{n^{k}}\to 0$ \ for every fixed $k>0$;
  \item $q^{n}\to 0$ \ if $|q|<1$; \quad $q^{n}\to\infty$ \ if $q>1$;
        \quad $q^{n}=1$ for $q=1$; \quad divergent for $q\leq -1$;
  \item $\sqrt[n]{n}\to 1$ \ and \ $\sqrt[n]{c}\to 1$ for every constant
        $c>0$;
  \item $\Bigl(1+\dfrac{x}{n}\Bigr)^{n}\to\mathrm{e}^{x}$ \ for every
        $x\in\R$; in particular
        $\Bigl(1+\dfrac{1}{n}\Bigr)^{n}\to\mathrm{e}\approx 2.71828$.
\end{itemize}
The last limit is the bridge to \emph{continuous compounding}: splitting a
yearly interest rate $x$ into $n$ compounding steps turns capital $K$ into
$K\,(1+x/n)^{n}$, and as the steps become infinitely fine this tends to
$K\,\mathrm{e}^{x}$, the formula behind continuous interest in financial
mathematics.
\end{rulebox}

\begin{rulebox}{Monotone convergence principle}
\textbf{Bounded $+$ monotone $\Rightarrow$ convergent.} An increasing
sequence that is bounded above converges; a decreasing sequence that is
bounded below converges. This is the workhorse for recursive sequences:
first establish monotonicity and boundedness, then find the candidate
limit $a$ from the \emph{fixed-point equation} obtained by letting
$n\to\infty$ in the recursion (e.g.\ $a=1.02\,a+100$ for
$a_{n+1}=1.02\,a_{n}+100$).
\end{rulebox}

\begin{rulebox}{Growth hierarchy}
For every $k>0$ and every $q>1$,
\[
  \ln n \;\ll\; n^{k} \;\ll\; q^{n},
  \qquad\text{meaning}\qquad
  \frac{\ln n}{n^{k}}\to 0
  \quad\text{and}\quad
  \frac{n^{k}}{q^{n}}\to 0 .
\]
Every power beats every logarithm, and every exponential with base $q>1$
beats every power. A useful consequence: $n^{k}q^{n}\to 0$ whenever
$|q|<1$, geometric decay swallows polynomial growth.
\end{rulebox}

\section{Practice Problems}

Compute the limit of each sequence or show that it diverges; where
indicated, also decide monotonicity and boundedness.

\begin{problems}
  \item $a_{n}=\dfrac{3n+5}{n+2}$
  \item $a_{n}=\dfrac{2n^{2}-n+4}{5n^{2}+3}$
  \item $a_{n}=\dfrac{n+1}{n^{2}+1}$
  \item $a_{n}=\dfrac{n^{3}-2n}{4n^{2}+n}$
  \item $a_{n}=\sqrt{n^{2}+n}-n$
  \item $a_{n}=\dfrac{\sqrt{4n^{2}+1}}{3n-2}$
  \item $a_{n}=(0.9)^{n}$ \ and \ $b_{n}=(1.1)^{n}$ \quad (limit or
        divergence for each)
  \item $a_{n}=\dfrac{2^{n}+5^{n}}{5^{n}}$
  \item $a_{n}=\dfrac{3^{n}-2^{n}}{3^{n}+2^{n}}$
  \item $a_{n}=\sqrt[n]{5n}$
  \item $a_{n}=\Bigl(1+\dfrac{3}{n}\Bigr)^{n}$
  \item $a_{n}=\Bigl(1-\dfrac{2}{n}\Bigr)^{n}$
  \item $a_{n}=\Bigl(1+\dfrac{1}{2n}\Bigr)^{n}$
  \item $a_{n}=\Bigl(\dfrac{n+4}{n}\Bigr)^{n}$
  \item $a_{n}=\dfrac{n}{n+1}$ \quad (monotone? bounded? convergent?,
        justify all three)
  \item $b_{n}=(-1)^{n}\,\dfrac{n}{n+1}$ \quad (monotone? bounded?
        convergent?)
  \item Savings account: $a_{1}=1000$ and $a_{n+1}=1.02\,a_{n}+100$
        (interest of $2\%$ per period plus a deposit of $100$). Compute
        $a_{2}$ and $a_{3}$, derive a closed formula for $a_{n}$, and
        describe the behaviour as $n\to\infty$.
  \item $a_{n}=\dfrac{2n+(-1)^{n}}{n}$ \quad (use the sandwich rule)
  \item $a_{n}=\dfrac{n^{2}}{2^{n}}$
  \item Classify as convergent or divergent, giving the limit where it
        exists:
        \[
          \text{(a)}\; a_{n}=(-1)^{n},\qquad
          \text{(b)}\; b_{n}=5-\frac{3}{n},\qquad
          \text{(c)}\; c_{n}=\ln n,\qquad
          \text{(d)}\; d_{n}=(-0.5)^{n}.
        \]
\end{problems}

\section{Solutions}

\begin{solutions}
  \item Divide numerator and denominator by $n$:
        \[
          a_{n}=\frac{3+\dfrac{5}{n}}{1+\dfrac{2}{n}}
          \;\longrightarrow\;\frac{3+0}{1+0}=3 ,
        \]
        using $5/n\to 0$, $2/n\to 0$ and the quotient rule. So
        $\lim_{n\to\infty}a_{n}=3$.

  \item Divide by the highest power $n^{2}$:
        \[
          a_{n}=\frac{2-\dfrac{1}{n}+\dfrac{4}{n^{2}}}{5+\dfrac{3}{n^{2}}}
          \;\longrightarrow\;\frac{2-0+0}{5+0}=\frac{2}{5} .
        \]

  \item Divide by $n^{2}$ (the highest power in the denominator):
        \[
          a_{n}=\frac{\dfrac{1}{n}+\dfrac{1}{n^{2}}}{1+\dfrac{1}{n^{2}}}
          \;\longrightarrow\;\frac{0+0}{1+0}=0 .
        \]
        Degree of the numerator below degree of the denominator: limit $0$.

  \item Divide by $n^{2}$:
        \[
          a_{n}=\frac{n-\dfrac{2}{n}}{4+\dfrac{1}{n}} .
        \]
        The denominator tends to $4$, while the numerator $n-2/n\to\infty$.
        Hence $a_{n}\to\infty$: the sequence is \textbf{divergent}
        (for large $n$, $a_{n}\approx n/4$).

  \item Multiply by the conjugate:
        \[
          a_{n}=\bigl(\sqrt{n^{2}+n}-n\bigr)
                \cdot\frac{\sqrt{n^{2}+n}+n}{\sqrt{n^{2}+n}+n}
               =\frac{(n^{2}+n)-n^{2}}{\sqrt{n^{2}+n}+n}
               =\frac{n}{\sqrt{n^{2}+n}+n} .
        \]
        Now divide numerator and denominator by $n$:
        \[
          a_{n}=\frac{1}{\sqrt{1+\dfrac{1}{n}}+1}
          \;\longrightarrow\;\frac{1}{\sqrt{1}+1}=\frac{1}{2} .
        \]

  \item Divide numerator and denominator by $n$ (inside the root: by
        $n^{2}$):
        \[
          a_{n}=\frac{\sqrt{4+\dfrac{1}{n^{2}}}}{3-\dfrac{2}{n}}
          \;\longrightarrow\;\frac{\sqrt{4}}{3}=\frac{2}{3} .
        \]

  \item For $a_{n}=(0.9)^{n}$: here $q=0.9$ with $|q|<1$, so
        $a_{n}\to 0$ (convergent). For $b_{n}=(1.1)^{n}$: here $q=1.1>1$,
        so $b_{n}\to\infty$ (divergent), capital growing at $10\%$ per
        period exceeds every bound.

  \item Split the fraction:
        \[
          a_{n}=\frac{2^{n}}{5^{n}}+\frac{5^{n}}{5^{n}}
               =\Bigl(\frac{2}{5}\Bigr)^{n}+1
          \;\longrightarrow\;0+1=1 ,
        \]
        since $|2/5|<1$.

  \item Divide numerator and denominator by the dominant term $3^{n}$:
        \[
          a_{n}=\frac{1-\bigl(\tfrac{2}{3}\bigr)^{n}}
                     {1+\bigl(\tfrac{2}{3}\bigr)^{n}}
          \;\longrightarrow\;\frac{1-0}{1+0}=1 ,
        \]
        because $(2/3)^{n}\to 0$ ($|q|<1$).

  \item Factor the root:
        \[
          a_{n}=\sqrt[n]{5n}=\sqrt[n]{5}\cdot\sqrt[n]{n}
          \;\longrightarrow\;1\cdot 1=1 ,
        \]
        by the standard limits $\sqrt[n]{c}\to 1$ ($c>0$) and
        $\sqrt[n]{n}\to 1$, combined with the product rule.

  \item This is the standard limit with $x=3$:
        \[
          \lim_{n\to\infty}\Bigl(1+\frac{3}{n}\Bigr)^{n}=\mathrm{e}^{3} .
        \]

  \item Standard limit with $x=-2$:
        \[
          \lim_{n\to\infty}\Bigl(1-\frac{2}{n}\Bigr)^{n}
          =\mathrm{e}^{-2}=\frac{1}{\mathrm{e}^{2}} .
        \]

  \item Rewrite the inner fraction so the pattern $x/n$ appears:
        $\dfrac{1}{2n}=\dfrac{1/2}{n}$, so
        \[
          a_{n}=\Bigl(1+\frac{1/2}{n}\Bigr)^{n}
          \;\longrightarrow\;\mathrm{e}^{1/2}=\sqrt{\mathrm{e}} .
        \]

  \item Simplify the base first:
        \[
          \frac{n+4}{n}=1+\frac{4}{n}
          \qquad\Longrightarrow\qquad
          a_{n}=\Bigl(1+\frac{4}{n}\Bigr)^{n}\to\mathrm{e}^{4} .
        \]

  \item \textbf{Monotone:} compute the difference
        \[
          a_{n+1}-a_{n}=\frac{n+1}{n+2}-\frac{n}{n+1}
          =\frac{(n+1)^{2}-n(n+2)}{(n+1)(n+2)}
          =\frac{n^{2}+2n+1-n^{2}-2n}{(n+1)(n+2)}
          =\frac{1}{(n+1)(n+2)}>0 ,
        \]
        so $(a_{n})$ is strictly increasing.
        \textbf{Bounded:} $0<a_{n}<1$ for all $n$, since $n<n+1$.
        \textbf{Convergent:} yes, by the monotone convergence principle
        (increasing and bounded above), and directly:
        \[
          a_{n}=\frac{1}{1+\dfrac{1}{n}}\;\longrightarrow\;\frac{1}{1+0}=1 .
        \]

  \item First terms: $b_{1}=-\tfrac{1}{2}$, $b_{2}=\tfrac{2}{3}$,
        $b_{3}=-\tfrac{3}{4}$, $b_{4}=\tfrac{4}{5}$. The signs alternate,
        so the sequence is \textbf{not monotone}. It is \textbf{bounded},
        because $|b_{n}|=\dfrac{n}{n+1}<1$ for all $n$. It is
        \textbf{divergent}: the even-indexed terms
        $b_{2m}=\tfrac{2m}{2m+1}\to 1$ while the odd-indexed terms tend to
        $-1$; a convergent sequence cannot approach two different values
        (uniqueness of the limit). Boundedness alone does not force
        convergence.

  \item First terms:
        \[
          a_{2}=1.02\cdot 1000+100=1020+100=1120 ,\qquad
          a_{3}=1.02\cdot 1120+100=1142.40+100=1242.40 .
        \]
        \textbf{Fixed point:} a constant solution $a^{*}$ would satisfy
        $a^{*}=1.02\,a^{*}+100$, i.e.\ $-0.02\,a^{*}=100$, so
        $a^{*}=-5000$. \textbf{Closed formula:} set
        $d_{n}=a_{n}-a^{*}=a_{n}+5000$. Then
        \[
          d_{n+1}=a_{n+1}+5000=1.02\,a_{n}+100+5000
                 =1.02\,(a_{n}+5000)=1.02\,d_{n} ,
        \]
        a geometric sequence with $d_{1}=1000+5000=6000$. Hence
        $d_{n}=6000\cdot(1.02)^{n-1}$ and
        \[
          a_{n}=6000\cdot(1.02)^{n-1}-5000 .
        \]
        Check: $a_{3}=6000\cdot 1.0404-5000=6242.40-5000=1242.40$,
        matches. \textbf{Behaviour:} $(a_{n})$ is strictly increasing
        ($a_{n+1}-a_{n}=0.02\,a_{n}+100>0$ since all terms are positive)
        and, because $1.02>1$ gives $(1.02)^{n-1}\to\infty$, the sequence
        \textbf{diverges to $\infty$}: interest plus steady deposits grow
        beyond every bound.

  \item Split: $a_{n}=2+\dfrac{(-1)^{n}}{n}$. Since
        $-1\leq(-1)^{n}\leq 1$,
        \[
          2-\frac{1}{n}\;\leq\;a_{n}\;\leq\;2+\frac{1}{n} .
        \]
        Both bounds tend to $2$, so by the sandwich rule
        $\lim_{n\to\infty}a_{n}=2$.

  \item Growth hierarchy: $n^{k}\ll q^{n}$ for $q>1$, here with $k=2$ and
        $q=2$, so
        \[
          \lim_{n\to\infty}\frac{n^{2}}{2^{n}}=0 .
        \]
        The decay is fast: for $n=10$ the term is
        $\tfrac{100}{1024}\approx 0.098$, for $n=20$ already
        $\tfrac{400}{1048576}\approx 0.0004$. Convergent with limit $0$.

  \item (a) $(-1)^{n}$ alternates between $-1$ and $1$ forever: bounded
        but \textbf{divergent} (no single limit). \
        (b) $3/n\to 0$, so $b_{n}\to 5-0=5$: \textbf{convergent} with
        limit $5$. \
        (c) $\ln n$ grows without bound (slowly, but beyond every $M$):
        \textbf{divergent} to $\infty$. \
        (d) $q=-0.5$ with $|q|<1$, so $d_{n}\to 0$: \textbf{convergent}
        with limit $0$.
\end{solutions}

\section{Additional Problems}

\begin{problems}
  \item $a_{n}=\dfrac{7n-4}{2n+9}$
  \item $a_{n}=\dfrac{n^{2}+6n}{3n^{2}-n+1}$
  \item $a_{n}=\dfrac{\sqrt{9n^{2}+n}}{n+1}$
  \item $a_{n}=\sqrt{n+1}-\sqrt{n}$
  \item $a_{n}=\dfrac{4^{n}-3^{n}}{4^{n}}$
  \item $a_{n}=\Bigl(1+\dfrac{5}{n}\Bigr)^{n}$
  \item $a_{n}=\Bigl(1-\dfrac{1}{n}\Bigr)^{2n}$
  \item $a_{n}=3+\dfrac{2}{n}$ \quad (monotone? bounded? convergent?,
        justify all three)
  \item Price adjustment: $p_{1}=200$ and $p_{n+1}=0.5\,p_{n}+60$ (each
        period the price moves halfway towards a demand-driven level plus
        a constant mark-up). Compute $p_{2}$ and $p_{3}$, derive a closed
        formula for $p_{n}$, and describe the behaviour as $n\to\infty$.
  \item Classify as convergent or divergent, giving the limit where it
        exists:
        \[
          \text{(a)}\; a_{n}=\frac{n}{2^{n}},\qquad
          \text{(b)}\; b_{n}=(1.05)^{n},\qquad
          \text{(c)}\; c_{n}=(-1)^{n}\,n,\qquad
          \text{(d)}\; d_{n}=\frac{2n+3}{n+1}.
        \]
\end{problems}

\section{Additional Solutions}

\begin{solutions}
  \item Divide numerator and denominator by $n$:
        \[
          a_{n}=\frac{7-\dfrac{4}{n}}{2+\dfrac{9}{n}}
          \;\longrightarrow\;\frac{7-0}{2+0}=\frac{7}{2} .
        \]

  \item Divide by $n^{2}$:
        \[
          a_{n}=\frac{1+\dfrac{6}{n}}{3-\dfrac{1}{n}+\dfrac{1}{n^{2}}}
          \;\longrightarrow\;\frac{1+0}{3-0+0}=\frac{1}{3} .
        \]

  \item Divide numerator and denominator by $n$; under the root this means
        dividing by $n^{2}$:
        \[
          a_{n}=\frac{\sqrt{9+\dfrac{1}{n}}}{1+\dfrac{1}{n}}
          \;\longrightarrow\;\frac{\sqrt{9}}{1}=3 .
        \]

  \item Multiply by the conjugate:
        \[
          a_{n}=\bigl(\sqrt{n+1}-\sqrt{n}\bigr)
                \cdot\frac{\sqrt{n+1}+\sqrt{n}}{\sqrt{n+1}+\sqrt{n}}
               =\frac{(n+1)-n}{\sqrt{n+1}+\sqrt{n}}
               =\frac{1}{\sqrt{n+1}+\sqrt{n}} .
        \]
        The denominator tends to $\infty$, hence $a_{n}\to 0$.

  \item Split the fraction:
        \[
          a_{n}=1-\Bigl(\frac{3}{4}\Bigr)^{n}
          \;\longrightarrow\;1-0=1 ,
        \]
        since $|3/4|<1$.

  \item Standard limit with $x=5$:
        \[
          \lim_{n\to\infty}\Bigl(1+\frac{5}{n}\Bigr)^{n}=\mathrm{e}^{5} .
        \]

  \item Write the power as a square of the standard form:
        \[
          a_{n}=\Bigl[\Bigl(1-\frac{1}{n}\Bigr)^{n}\Bigr]^{2}
          \;\longrightarrow\;\bigl(\mathrm{e}^{-1}\bigr)^{2}
          =\mathrm{e}^{-2}=\frac{1}{\mathrm{e}^{2}} ,
        \]
        using $(1-1/n)^{n}\to\mathrm{e}^{-1}$ and the product rule.

  \item \textbf{Monotone:}
        \[
          a_{n+1}-a_{n}=\frac{2}{n+1}-\frac{2}{n}
          =\frac{2n-2(n+1)}{n(n+1)}=\frac{-2}{n(n+1)}<0 ,
        \]
        so $(a_{n})$ is strictly decreasing.
        \textbf{Bounded:} $3<a_{n}\leq a_{1}=5$ for all $n$.
        \textbf{Convergent:} yes: decreasing and bounded below, so the
        monotone convergence principle applies; since $2/n\to 0$, the
        limit is $3$.

  \item First terms:
        \[
          p_{2}=0.5\cdot 200+60=100+60=160 ,\qquad
          p_{3}=0.5\cdot 160+60=80+60=140 .
        \]
        \textbf{Fixed point:} $p^{*}=0.5\,p^{*}+60$ gives
        $0.5\,p^{*}=60$, so $p^{*}=120$. \textbf{Closed formula:} set
        $d_{n}=p_{n}-120$. Then
        \[
          d_{n+1}=p_{n+1}-120=0.5\,p_{n}+60-120=0.5\,(p_{n}-120)
                 =0.5\,d_{n} ,
        \]
        with $d_{1}=200-120=80$, so $d_{n}=80\cdot(0.5)^{n-1}$ and
        \[
          p_{n}=120+80\cdot(0.5)^{n-1} .
        \]
        Check: $p_{3}=120+80\cdot 0.25=120+20=140$, matches.
        \textbf{Behaviour:} $(0.5)^{n-1}\to 0$ because $|0.5|<1$, so
        $p_{n}\to 120$: the price falls monotonically to the equilibrium
        $120$. Contrast with the savings account of Problem~17 in the main
        list: a coefficient below $1$ pulls the sequence \emph{towards}
        the fixed point, a coefficient above $1$ (there $1.02$) pushes it
        \emph{away}.

  \item (a) Growth hierarchy ($n^{k}\ll 2^{n}$ with $k=1$):
        $a_{n}\to 0$, \textbf{convergent}. \
        (b) $q=1.05>1$, so $b_{n}\to\infty$: \textbf{divergent}. \
        (c) $|c_{n}|=n\to\infty$ with alternating signs: unbounded, hence
        \textbf{divergent} (no limit, not even $\pm\infty$). \
        (d) Divide by $n$:
        $d_{n}=\dfrac{2+3/n}{1+1/n}\to\dfrac{2}{1}=2$:
        \textbf{convergent} with limit $2$.
\end{solutions}

\chapter{Series}

A series adds up the terms of a sequence. Since infinitely many
additions cannot be carried out one after another, the value of a series is
\emph{defined} as the limit of its partial sums: series theory is sequence
theory applied to one particular sequence. This chapter covers the
convergence basics: the two sum formulas everyone must know (arithmetic and
geometric), the necessary condition for convergence with the harmonic series
as its famous counterexample, telescoping sums, and two light convergence
tests. The star of the chapter is the geometric series: the entire financial
mathematics of the next chapter (compound interest, present value,
annuities) is the geometric series in a suit.

\section{Core Rules}

\begin{rulebox}{Partial sums and convergence}
Given a sequence $(a_k)_{k\ge 1}$, its $n$-th \emph{partial sum} is
\[
  s_n \;=\; \sum_{k=1}^{n} a_k \;=\; a_1+a_2+\dots+a_n .
\]
The series $\sum_{k=1}^{\infty} a_k$ \emph{converges} if the sequence
$(s_n)$ of partial sums converges, and then
\[
  \sum_{k=1}^{\infty} a_k \;=\; \lim_{n\to\infty} s_n .
\]
If $(s_n)$ has no limit, the series \emph{diverges}. A series is thus
nothing but a sequence in disguise: every statement about series is a
statement about its sequence of partial sums.
\end{rulebox}

\begin{rulebox}{Arithmetic series (finite sum)}
For an arithmetic sequence $a_k=a+(k-1)d$ the sum of the first $n$ terms is
\[
  \sum_{k=1}^{n}\bigl(a+(k-1)d\bigr)
  \;=\; \frac{n}{2}\bigl(2a+(n-1)d\bigr)
  \;=\; n\cdot\frac{a_1+a_n}{2}
\],
$n$ times the average of the first and the last term. The special case
$a=d=1$ is Gauss' formula
\[
  \sum_{k=1}^{n} k \;=\; \frac{n(n+1)}{2} .
\]
An \emph{infinite} arithmetic series never
converges unless all terms are $0$: its terms do not tend to $0$.
\end{rulebox}

\begin{rulebox}{Geometric series, finite sum}
For a geometric sequence $a q^{k}$ with ratio $q\neq 1$,
\[
  \sum_{k=0}^{n-1} a q^{k}
  \;=\; a\,\frac{1-q^{n}}{1-q} \qquad (n \text{ terms}).
\]
For $q=1$ the sum is simply $na$.
\end{rulebox}

\begin{rulebox}{Geometric series, infinite sum}
If $|q|<1$, then $q^{n}\to 0$, so the partial sums converge and
\[
  \sum_{k=0}^{\infty} a q^{k} \;=\; \frac{a}{1-q},
  \qquad\text{and starting at } k=1:\quad
  \sum_{k=1}^{\infty} a q^{k} \;=\; \frac{aq}{1-q} .
\]
For $|q|\ge 1$ the series diverges. This is \textbf{the} key formula of the
whole book: in the next chapter every present value, annuity and perpetuity
formula of financial mathematics is obtained by choosing $q$ as a discount
factor in the geometric series.
\end{rulebox}

\begin{rulebox}{Necessary condition: the terms must vanish}
If $\sum_{k=1}^{\infty} a_k$ converges, then $a_k\to 0$. In contrapositive
form this is the everyday \emph{divergence test}:
\[
  a_k \not\to 0 \quad\Longrightarrow\quad \sum_{k=1}^{\infty} a_k
  \text{ diverges.}
\]
The converse is \textbf{false}: $a_k\to 0$ does \emph{not} guarantee
convergence. The standard counterexample is the harmonic series
\[
  \sum_{k=1}^{\infty}\frac{1}{k} \;=\; 1+\frac12+\frac13+\frac14+\dots,
\]
whose terms tend to $0$ while the partial sums grow beyond every bound.
\end{rulebox}

\begin{rulebox}{Telescoping sums}
If the terms can be written as consecutive differences, $a_k=b_k-b_{k+1}$,
then everything in the middle cancels:
\[
  \sum_{k=1}^{n}(b_k-b_{k+1}) \;=\; b_1-b_{n+1} .
\]
The series converges precisely when $(b_n)$ converges, with sum
$b_1-\lim_{n\to\infty} b_n$. A telescoping sum is
usually uncovered by partial fractions, e.g.
$\dfrac{1}{k(k+1)}=\dfrac{1}{k}-\dfrac{1}{k+1}$.
\end{rulebox}

\begin{rulebox}{Comparison test (simple form)}
Suppose $0\le a_k\le b_k$ for all $k$ (or from some index on).
\begin{itemize}[leftmargin=1.6em,itemsep=0.3ex]
  \item \textbf{Majorant}: if
        $\sum b_k$ converges, so does $\sum a_k$.
  \item \textbf{Minorant}: if
        $\sum a_k$ diverges, so does $\sum b_k$.
\end{itemize}
The usual benchmarks are the geometric series (convergent for $|q|<1$) and
the harmonic series (divergent). Example:
$\dfrac{1}{2^{k}+1}\le\Bigl(\dfrac12\Bigr)^{k}$, so
$\sum\dfrac{1}{2^{k}+1}$ converges.
\end{rulebox}

\begin{rulebox}{Ratio test (light form)}
For $a_k>0$ suppose the ratio of consecutive terms has a limit,
$\dfrac{a_{k+1}}{a_k}\to L$. Then:
\begin{itemize}[leftmargin=1.6em,itemsep=0.3ex]
  \item $L<1$: the series converges (it behaves like a geometric series
        with ratio $L$);
  \item $L>1$: the series diverges (the terms eventually grow);
  \item $L=1$: no conclusion: both $\sum 1/k$ (divergent) and
        $\sum 1/k^{2}$ (convergent) give $L=1$.
\end{itemize}
Example: for $a_k=k/2^{k}$ the ratio is
$\dfrac{k+1}{2k}\to\dfrac12<1$, so $\sum k/2^{k}$ converges. The test shines when powers or factorials appear.
\end{rulebox}

\begin{rulebox}{Linearity}
If $\sum a_k=A$ and $\sum b_k=B$ both converge and $\lambda,\mu\in\R$, then
\[
  \sum_{k=1}^{\infty}\bigl(\lambda a_k+\mu b_k\bigr)
  \;=\; \lambda A+\mu B .
\]
Splitting a series into pieces is only legitimate when \emph{each} piece
converges on its own.
\end{rulebox}

\section{Practice Problems}

Evaluate the sums in Problems 1--14. In Problems 15--19 decide convergence
or divergence with a justified test. Problem 20 previews the next chapter.

\begin{problems}
  \item $\displaystyle\sum_{k=1}^{100} k$
  \item $7+10+13+\dots+118$
  \item $\displaystyle\sum_{k=1}^{50}(2k-1)$ \quad (the first $50$ odd
        numbers)
  \item $\displaystyle\sum_{k=0}^{9} 2^{k}$
  \item $\displaystyle\sum_{k=1}^{6} 5\cdot 3^{k-1}$
  \item $\displaystyle\sum_{k=0}^{5}\Bigl(\frac12\Bigr)^{k}$
  \item $\displaystyle\sum_{k=0}^{\infty}\Bigl(\frac13\Bigr)^{k}$
  \item $\displaystyle\sum_{k=1}^{\infty}\Bigl(\frac25\Bigr)^{k}$
  \item $\displaystyle\sum_{k=0}^{\infty}\Bigl(-\frac12\Bigr)^{k}$
  \item Show that $0.999\ldots=0.\overline{9}$ equals exactly $1$.
  \item Write $2.353535\ldots=2.\overline{35}$ as a fraction.
  \item $\displaystyle\sum_{k=1}^{n}\frac{1}{k(k+1)}$ \quad (find a closed
        formula, then the value of the infinite series)
  \item $\displaystyle\sum_{k=1}^{99}\bigl(\sqrt{k+1}-\sqrt{k}\bigr)$
  \item $\displaystyle\sum_{k=1}^{\infty}\frac{1}{(2k-1)(2k+1)}$
  \item Does $\displaystyle\sum_{k=1}^{\infty}\frac{k}{k+1}$ converge?
  \item The terms of $\displaystyle\sum_{k=1}^{\infty}\frac{1}{k}$ satisfy
        $\frac{1}{k}\to 0$. Does the series converge? Justify carefully.
  \item Does $\displaystyle\sum_{k=1}^{\infty}\frac{1}{2^{k}+1}$ converge?
  \item Does $\displaystyle\sum_{k=1}^{\infty}\frac{1}{\sqrt{k}}$ converge?
  \item Does $\displaystyle\sum_{k=1}^{\infty}\frac{k}{2^{k}}$ converge?
  \item \emph{Perpetuity teaser.} A contract pays $100$ euros at the end of
        every year, forever. At an interest rate of $5\%$ per year, a
        payment $k$ years from now is worth $100\cdot(1/1.05)^{k}$ euros
        today. What is the total value of the contract today?
\end{problems}

\section{Solutions}

\begin{solutions}
  \item Arithmetic with $a_1=1$, $a_{100}=100$, $n=100$:
        \[
          \sum_{k=1}^{100} k=\frac{100\cdot 101}{2}=5050 .
        \]

  \item Arithmetic with $a=7$, $d=3$. The last term fixes $n$:
        $118=7+3(n-1)$ gives $3(n-1)=111$, so $n=38$. Hence
        \[
          7+10+\dots+118=\frac{38}{2}\,(7+118)=19\cdot 125=2375 .
        \]

  \item Arithmetic with $a_1=1$, $a_{50}=99$, $n=50$:
        \[
          \sum_{k=1}^{50}(2k-1)=50\cdot\frac{1+99}{2}=50\cdot 50=2500 .
        \]
        In general the first $n$ odd numbers sum to $n^{2}$.

  \item Geometric, $a=1$, $q=2$, $10$ terms:
        \[
          \sum_{k=0}^{9}2^{k}=\frac{1-2^{10}}{1-2}=\frac{1-1024}{-1}=1023 .
        \]

  \item Pull out the factor $5$; the remaining sum is geometric with $q=3$
        and $6$ terms:
        \[
          \sum_{k=1}^{6}5\cdot 3^{k-1}
          =5\sum_{j=0}^{5}3^{j}
          =5\cdot\frac{1-3^{6}}{1-3}
          =5\cdot\frac{-728}{-2}=5\cdot 364=1820 .
        \]

  \item Geometric, $q=\tfrac12$, $6$ terms:
        \[
          \sum_{k=0}^{5}\Bigl(\frac12\Bigr)^{k}
          =\frac{1-(1/2)^{6}}{1-1/2}
          =2\Bigl(1-\frac{1}{64}\Bigr)=\frac{63}{32} .
        \]

  \item $|q|=\tfrac13<1$, so the infinite sum formula applies:
        \[
          \sum_{k=0}^{\infty}\Bigl(\frac13\Bigr)^{k}
          =\frac{1}{1-\tfrac13}=\frac{3}{2} .
        \]

  \item The sum starts at $k=1$, so the first term is $\tfrac25$ and
        $q=\tfrac25$:
        \[
          \sum_{k=1}^{\infty}\Bigl(\frac25\Bigr)^{k}
          =\frac{2/5}{1-2/5}=\frac{2/5}{3/5}=\frac{2}{3} .
        \]

  \item The formula also covers negative ratios; here $q=-\tfrac12$ with
        $|q|<1$:
        \[
          \sum_{k=0}^{\infty}\Bigl(-\frac12\Bigr)^{k}
          =\frac{1}{1-(-\tfrac12)}=\frac{1}{3/2}=\frac{2}{3} .
        \]

  \item The decimal expansion is a geometric series with $q=\tfrac{1}{10}$:
        \[
          0.\overline{9}
          =\sum_{k=1}^{\infty}\frac{9}{10^{k}}
          =\frac{9\cdot\tfrac{1}{10}}{1-\tfrac{1}{10}}
          =\frac{9/10}{9/10}=1 .
        \]
        The equality is exact: $0.\overline{9}$ is not ``a little less
        than $1$'', it \emph{is} $1$.

  \item Split off the integer part; the repeating block $35$ has period
        $2$, so $q=\tfrac{1}{100}$:
        \[
          2.\overline{35}
          =2+\sum_{k=1}^{\infty}\frac{35}{100^{k}}
          =2+\frac{35\cdot\tfrac{1}{100}}{1-\tfrac{1}{100}}
          =2+\frac{35}{99}=\frac{233}{99} .
        \]

  \item Partial fractions: $\dfrac{1}{k(k+1)}=\dfrac{1}{k}-\dfrac{1}{k+1}$,
        so the sum telescopes:
        \[
          s_n=\sum_{k=1}^{n}\Bigl(\frac{1}{k}-\frac{1}{k+1}\Bigr)
          =1-\frac{1}{n+1} .
        \]
        Since $\dfrac{1}{n+1}\to 0$,
        \[
          \sum_{k=1}^{\infty}\frac{1}{k(k+1)}=\lim_{n\to\infty} s_n=1 .
        \]

  \item Telescoping with $b_k=-\sqrt{k}$ (or simply: all middle terms
        cancel):
        \[
          \sum_{k=1}^{99}\bigl(\sqrt{k+1}-\sqrt{k}\bigr)
          =\sqrt{100}-\sqrt{1}=10-1=9 .
        \]

  \item Partial fractions:
        \[
          \frac{1}{(2k-1)(2k+1)}
          =\frac12\Bigl(\frac{1}{2k-1}-\frac{1}{2k+1}\Bigr).
        \]
        The partial sums telescope:
        \[
          s_n=\frac12\Bigl(1-\frac{1}{2n+1}\Bigr)
          \;\longrightarrow\;\frac12 ,
        \]
        so the series converges with sum $\dfrac12$.

  \item No. The terms satisfy
        \[
          a_k=\frac{k}{k+1}\;\longrightarrow\;1\neq 0,
        \]
        so the necessary condition fails and the series diverges. (Each
        late term adds almost $1$, so the partial sums grow roughly like
        $n$.)

  \item No: this is the harmonic series, the standard warning that
        $a_k\to 0$ is necessary but \emph{not sufficient}. Group the terms
        in blocks of doubling length:
        \[
          1+\frac12
          +\underbrace{\Bigl(\frac13+\frac14\Bigr)}_{\ge\,2\cdot\frac14=\frac12}
          +\underbrace{\Bigl(\frac15+\frac16+\frac17+\frac18\Bigr)}_{\ge\,4\cdot\frac18=\frac12}
          +\dots
        \]
        Each block contributes at least $\tfrac12$, so
        $s_{2^{m}}\ge 1+\tfrac{m}{2}$. The partial sums exceed every bound,
        hence the series diverges.

  \item Yes, by comparison with a geometric majorant: for all $k\ge 1$,
        \[
          0<\frac{1}{2^{k}+1}<\Bigl(\frac12\Bigr)^{k},
          \qquad
          \sum_{k=1}^{\infty}\Bigl(\frac12\Bigr)^{k}
          =\frac{1/2}{1-1/2}=1<\infty ,
        \]
        so $\displaystyle\sum_{k=1}^{\infty}\frac{1}{2^{k}+1}$ converges.

  \item No. For $k\ge 1$ we have $\sqrt{k}\le k$, hence
        \[
          \frac{1}{\sqrt{k}}\ \ge\ \frac{1}{k} .
        \]
        The harmonic series diverges, so by the minorant criterion
        $\displaystyle\sum_{k=1}^{\infty}\frac{1}{\sqrt{k}}$ diverges as
        well.

  \item Yes, by the ratio test with $a_k=k/2^{k}$:
        \[
          \frac{a_{k+1}}{a_k}
          =\frac{k+1}{2^{k+1}}\cdot\frac{2^{k}}{k}
          =\frac{k+1}{2k}
          \;\longrightarrow\;\frac12<1 ,
        \]
        so the series converges. (Its value happens to be $2$, but the
        test alone does not reveal that.)

  \item The total present value is a geometric series with
        $q=\dfrac{1}{1.05}$, starting at $k=1$:
        \[
          \mathrm{PV}
          =\sum_{k=1}^{\infty}100\,\Bigl(\frac{1}{1.05}\Bigr)^{k}
          =100\cdot\frac{\tfrac{1}{1.05}}{1-\tfrac{1}{1.05}}
          =100\cdot\frac{\tfrac{1}{1.05}}{\tfrac{0.05}{1.05}}
          =\frac{100}{0.05}=2000 .
        \]
        An infinite stream of payments is worth a finite amount:
        $2000$ euros. The rule of thumb
        $\mathrm{PV}=\text{payment}/\text{interest rate}$ is exactly the
        infinite geometric sum, the next chapter builds all of financial
        mathematics on this observation.
\end{solutions}

\section{Additional Problems}

\begin{problems}
  \item $\displaystyle\sum_{k=1}^{200}(3k+2)$
  \item $100+96+92+\dots+4$
  \item $\displaystyle\sum_{k=0}^{7}3\cdot(-2)^{k}$
  \item $\displaystyle\sum_{k=2}^{\infty}\Bigl(\frac14\Bigr)^{k}$
  \item Write $0.727272\ldots=0.\overline{72}$ as a fraction.
  \item $\displaystyle\sum_{k=1}^{\infty}\frac{1}{(k+1)(k+2)}$
  \item Does $\displaystyle\sum_{k=1}^{\infty}\frac{2k+1}{3k+5}$ converge?
  \item Show that $\displaystyle\sum_{k=1}^{\infty}\frac{1}{k^{2}}$
        converges, by comparing with the telescoping series of
        Problem~12 of the practice set
        \emph{(hint: $\tfrac{1}{k^{2}}\le\tfrac{2}{k(k+1)}$)}.
  \item Does $\displaystyle\sum_{k=1}^{\infty}\frac{2^{k}}{k!}$ converge?
  \item $\displaystyle\sum_{k=0}^{\infty}
        \Bigl(\Bigl(\frac12\Bigr)^{k}+2\cdot\Bigl(\frac13\Bigr)^{k}\Bigr)$
\end{problems}

\section{Additional Solutions}

\begin{solutions}
  \item Split by linearity of finite sums and use Gauss' formula:
        \[
          \sum_{k=1}^{200}(3k+2)
          =3\sum_{k=1}^{200}k+\sum_{k=1}^{200}2
          =3\cdot\frac{200\cdot 201}{2}+400
          =3\cdot 20100+400=60700 .
        \]

  \item Arithmetic with $a=100$, $d=-4$. The last term fixes $n$:
        $4=100-4(n-1)$ gives $4(n-1)=96$, so $n=25$. Hence
        \[
          100+96+\dots+4=\frac{25}{2}\,(100+4)=25\cdot 52=1300 .
        \]

  \item Geometric with $q=-2$, $8$ terms:
        \[
          \sum_{k=0}^{7}3\cdot(-2)^{k}
          =3\cdot\frac{1-(-2)^{8}}{1-(-2)}
          =3\cdot\frac{1-256}{3}=-255 .
        \]

  \item Geometric with $|q|=\tfrac14<1$; the first term is $(1/4)^{2}$:
        \[
          \sum_{k=2}^{\infty}\Bigl(\frac14\Bigr)^{k}
          =\frac{(1/4)^{2}}{1-1/4}
          =\frac{1/16}{3/4}=\frac{1}{12} .
        \]
        (Check: full sum $\tfrac{4}{3}$ minus the terms $1+\tfrac14$
        indeed leaves $\tfrac{4}{3}-\tfrac{5}{4}=\tfrac{1}{12}$.)

  \item Repeating block $72$ of period $2$, so $q=\tfrac{1}{100}$:
        \[
          0.\overline{72}
          =\sum_{k=1}^{\infty}\frac{72}{100^{k}}
          =\frac{72\cdot\tfrac{1}{100}}{1-\tfrac{1}{100}}
          =\frac{72}{99}=\frac{8}{11} .
        \]

  \item Partial fractions:
        $\dfrac{1}{(k+1)(k+2)}=\dfrac{1}{k+1}-\dfrac{1}{k+2}$, so
        \[
          s_n=\frac12-\frac{1}{n+2}\;\longrightarrow\;\frac12 ,
        \]
        and the series converges with sum $\dfrac12$.

  \item No. The terms satisfy
        \[
          \frac{2k+1}{3k+5}
          =\frac{2+\tfrac{1}{k}}{3+\tfrac{5}{k}}
          \;\longrightarrow\;\frac{2}{3}\neq 0 ,
        \]
        so the necessary condition fails and the series diverges.

  \item For $k\ge 1$ we have $k+1\le 2k$, hence $k(k+1)\le 2k^{2}$ and
        therefore
        \[
          \frac{1}{k^{2}}\ \le\ \frac{2}{k(k+1)} .
        \]
        The majorant is twice the telescoping series with sum $1$, so
        $\sum_{k=1}^{\infty}\tfrac{2}{k(k+1)}=2$ converges. By the
        comparison test, $\displaystyle\sum_{k=1}^{\infty}\frac{1}{k^{2}}$
        converges (its partial sums are increasing and bounded by $2$).
        Its exact value, $\pi^{2}/6$, is a famous result of Euler: far
        beyond what the test tells us.

  \item Yes, by the ratio test with $a_k=2^{k}/k!$:
        \[
          \frac{a_{k+1}}{a_k}
          =\frac{2^{k+1}}{(k+1)!}\cdot\frac{k!}{2^{k}}
          =\frac{2}{k+1}
          \;\longrightarrow\;0<1 ,
        \]
        so the series converges. (Its value is $\mathrm{e}^{2}-1$, as the
        exponential series will show later.)

  \item Both pieces are convergent geometric series, so linearity applies:
        \[
          \sum_{k=0}^{\infty}\Bigl(\frac12\Bigr)^{k}
          =\frac{1}{1-\tfrac12}=2 ,
          \qquad
          2\sum_{k=0}^{\infty}\Bigl(\frac13\Bigr)^{k}
          =2\cdot\frac{1}{1-\tfrac13}=2\cdot\frac32=3 ,
        \]
        hence the total is $2+3=5$.
\end{solutions}

\chapter{Financial Mathematics}

Everything in this chapter is an application of the geometric series. A
capital under compound interest is a geometric sequence; a stream of equal
payments, accumulated or discounted, is a geometric sum; and the central
formula of the previous chapter,
\[
  \sum_{k=0}^{n-1}q^{k}=\frac{q^{n}-1}{q-1}\qquad(q\neq 1),
\]
turns directly into the end value of a savings plan, the present value of
an annuity, the payment on a loan and (via the limit $n\to\infty$)
the value of a perpetuity.

Time is discrete throughout: interest is credited at the end of each
period, and payments arrive at period ends (or period beginnings, for the
annuity-due variant). Continuous compounding $K_{0}\mathrm{e}^{rt}$
appears exactly once, as a limit that bridges to later volumes;
continuous \emph{discounting} of payment streams by integrals is drilled
in \chref{integrals}{7}. All money amounts are in
euros.

\section{Core Rules}

\begin{rulebox}{Simple vs.\ compound interest}
Let $K_{0}$ be the initial capital and $i$ the annual interest rate
(p.a., written as a decimal). Under \emph{simple interest} only $K_{0}$ earns interest:
\[
  K_{n}=K_{0}(1+n\,i) .
\]
Under \emph{compound interest} credited interest itself
earns interest. With the accumulation factor $q=1+i$,
\[
  K_{n}=K_{0}\,q^{n}=K_{0}(1+i)^{n} .
\]
The capitals $K_{0},K_{1},K_{2},\dots$ form a geometric sequence with
ratio $q$; simple interest grows linearly, compound interest
exponentially.
\end{rulebox}

\begin{rulebox}{Sub-annual compounding and the effective annual rate}
If a nominal rate $j$ p.a.\ is credited in $m$ equal
instalments of $j/m$ per period ($m=4$: quarterly, $m=12$: monthly), then
after $t$ years
\[
  K=K_{0}\left(1+\frac{j}{m}\right)^{mt},
\]
and the \emph{effective annual rate} is the
annually compounded rate producing the same growth:
\[
  i_{\mathrm{eff}}=\left(1+\frac{j}{m}\right)^{m}-1 .
\]
For $m>1$ compounding within the year makes
$i_{\mathrm{eff}}>j$. Competing offers are compared by their effective
rates, never by their nominal rates.
\end{rulebox}

\begin{rulebox}{Continuous compounding (bridge)}
Letting the number of instalments grow without bound and using the limit
$\left(1+\frac{r}{m}\right)^{m}\to\mathrm{e}^{r}$ from the chapter on
sequences,
\[
  K(t)=\lim_{m\to\infty}K_{0}\left(1+\frac{r}{m}\right)^{mt}
      =K_{0}\,\mathrm{e}^{rt} .
\]
This is the largest end value a nominal rate $r$ can generate, and the
form in which interest enters most theoretical models. In this chapter it
serves only as this limit; continuous discounting of entire payment
streams (an integral, not a sum) is drilled in \chref{integrals}{7}.
\end{rulebox}

\begin{rulebox}{Discounting and present value}
Solving $K_{n}=K_{0}q^{n}$ for $K_{0}$ gives the \emph{present value} of a payment $K_{n}$ due in $n$ years:
\[
  K_{0}=\frac{K_{n}}{q^{n}}=K_{n}(1+i)^{-n} .
\]
The factor $q^{-1}=\frac{1}{1+i}$ is the discount factor. Payments due at different dates are made
comparable by discounting all of them to today, this single idea
underlies annuities, loans and project appraisal below.
\end{rulebox}

\begin{rulebox}{Annuity: end value and present value}
An ordinary annuity pays the constant amount
$R$ at the \emph{end} of each of the years $1,\dots,n$. Immediately after
the last payment, payment $k$ has earned interest for $n-k$ years, so the
end value is a geometric sum:
\[
  E_{n}=R\sum_{k=1}^{n}q^{\,n-k}=R\sum_{j=0}^{n-1}q^{\,j}
       =R\,\frac{q^{n}-1}{q-1}=R\,\frac{q^{n}-1}{i} .
\]
Discounting $E_{n}$ back $n$ years gives the present value
\[
  B_{n}=E_{n}\,q^{-n}=R\,\frac{1-q^{-n}}{i} .
\]
For an annuity-due, payments at the
\emph{beginnings} of the years, every payment earns one extra year of
interest: multiply $E_{n}$ and $B_{n}$ by $q$.
\end{rulebox}

\begin{rulebox}{Perpetuity}
Let $n\to\infty$ in the present value $B_{n}=R\,\frac{1-q^{-n}}{i}$.
Since $q>1$ we have $q^{-n}\to 0$ (geometric decay, as in the previous
chapter), so a perpetual annual payment $R$ is worth
\[
  B_{\infty}=\frac{R}{i}
\]
today. Read in both directions: a capital $C$ sustains the perpetual
payment $R=C\,i$, and a perpetual payment $R$ bought for the price $P$
implies the rate $i=R/P$.
\end{rulebox}

\begin{rulebox}{Loan amortization}
A debt $S$ at rate $i$ is repaid by $n$ equal annual payments $A$ if the present value of the payments equals
the debt: $S=A\,\frac{1-q^{-n}}{i}$, hence
\[
  A=S\,\frac{i}{1-q^{-n}}=S\,\frac{i\,q^{n}}{q^{n}-1} .
\]
Each payment splits into interest, $i$ times the debt still outstanding,
and principal, which reduces the debt; as the debt falls,
the interest share shrinks and the principal share grows.
Example: $S=1000$, $i=10\%$, $n=2$:
$A=1000\cdot\frac{0.1\cdot 1.21}{0.21}=576.19$, and the schedule closes
exactly:

\begin{center}
\begin{tabular}{c r r r r}
\hline
Year & Debt at start & Interest & Principal & Debt at end\\
\hline
$1$ & $1000.00$ & $100.00$ & $476.19$ & $523.81$\\
$2$ & $523.81$ & $52.38$ & $523.81$ & $0.00$\\
\hline
\end{tabular}
\end{center}
\end{rulebox}

\begin{rulebox}{Net present value}
A project paying the cash flows $C_{0},C_{1},\dots,C_{n}$ at the ends of
years $0,1,\dots,n$ (typically $C_{0}<0$, today's outlay) has
\emph{net present value}
\[
  \mathrm{NPV}=\sum_{t=0}^{n}\frac{C_{t}}{q^{t}}
\]
at the discount rate $i$ (the opportunity cost of capital). Decision
rule: invest if $\mathrm{NPV}>0$; among mutually exclusive projects
evaluated at the same rate, choose the largest NPV. Undiscounted totals
are meaningless: a euro in year~2 is not a euro today.
\end{rulebox}

\section{Worked Example}

\begin{worked}{A $200{,}000$ EUR mortgage at $4\%$ nominal, compounded monthly, over $20$ years: find the monthly payment and the effective annual rate; then appraise a project paying $400$ for three years against an outlay of $1000$ at $10\%$}
\wstep{Step 1. Convert the quoted rate to a \emph{per-period} rate.}
``$4\%$ nominal, compounded monthly'' means $4\%$ divided into twelve equal
pieces:
\[
  i=\frac{0.04}{12}=0.00\overline{3}\ \text{per month},
  \qquad
  n=20\times12=240\ \text{periods}.
\]
\wtrap{The rate and the number of periods must refer to the \emph{same}
unit of time. Using $i=0.04$ with $n=240$, or $i=0.00\overline{3}$ with
$n=20$, are the two ways this goes wrong, and both produce answers that look
plausible. Write the units next to both numbers.}

\wstep{Step 2. Recognise the cash-flow shape: an ordinary annuity.}
The bank lends $P$ today against $240$ equal payments $R$ at the \emph{end}
of each month. Present value of the payments $=$ amount lent:
\[
  P=R\cdot\frac{1-(1+i)^{-n}}{i}
  \qquad\Longrightarrow\qquad
  R=P\cdot\frac{i}{1-(1+i)^{-n}} .
\]

\wstep{Step 3. Evaluate the discount factor, carefully.}
\[
  (1+i)^{-240}=(1.00\overline{3})^{-240}=0.449927 ,
  \qquad
  1-(1+i)^{-n}=0.550073 .
\]
\[
  R=200\,000\cdot\frac{0.00\overline{3}}{0.550073}
   =\frac{666.667}{0.550073}
   =\boxed{1{,}211.96\ \text{EUR per month}}
\]
\wtrap{Keep six significant figures until the last line. Rounding
$(1+i)^{-240}$ to $0.45$ shifts the payment by about a euro a month --- 
which over $240$ months is a few hundred euros, and in an examination is a
wrong answer.}

\wstep{Step 4. Check by running the formula backwards.}
$1211.96\times\frac{1-(1.00\overline{3})^{-240}}{0.00\overline{3}}
=199\,999.9$, which is $200{,}000$ EUR up to the rounding of the payment.
\checkmark\ Total paid: $240\times1211.96=290{,}871$ EUR, so
$90{,}871$ EUR is interest --- a useful reality check on the size of the
answer.

\wstep{Step 5. The effective annual rate: compounding is not addition.}
Twelve months of $0.00\overline{3}$ do not give $4\%$, because interest
earns interest:
\[
  \mathrm{EAR}=\Bigl(1+\frac{0.04}{12}\Bigr)^{12}-1=1.040742-1=\boxed{4.0742\%}
\]
The gap of $7.4$ basis points is the compounding effect. Continuous
compounding is the limiting case, $\mathrm{e}^{0.04}-1=4.0811\%$, and it is
an upper bound: no compounding frequency can exceed it.
\wtrap{Compare loans by EAR, never by nominal rate. And do not subtract the
$1$ too early --- the ``$-1$'' converts a growth \emph{factor} into a
\emph{rate}, and omitting it turns $4\%$ into $104\%$.}

\wstep{Step 6. The project: discount every flow to today.}
Outlay $1000$ now, then $400$ at the end of each of three years, at $r=10\%$:
\[
  \mathrm{NPV}=-1000+\frac{400}{1.1}+\frac{400}{1.1^{2}}+\frac{400}{1.1^{3}}
  =-1000+363.64+330.58+300.53=\boxed{-5.26\ \text{EUR}}
\]
Negative, so \textbf{reject}. Note how close to zero it is: the project
almost pays for itself, which is why inspection alone cannot decide here
and the calculation is needed.
\wtrap{The first flow is at $t=0$ and is \emph{not} discounted --- exponent
$0$, not $1$. Discounting the outlay as well is the most common error in
NPV, and it flatters every project.}

\wstep{Step 7. Cross-check with the internal rate of return.}
The IRR solves $\mathrm{NPV}(r)=0$; numerically $r^{*}\approx9.70\%$. Since
$9.70\%<10\%$, the project earns less than the required return --- the same
verdict as the NPV, reached independently.
\[
  \mathrm{NPV}>0 \iff \mathrm{IRR}>r
  \quad\text{(for a conventional cash flow: one sign change).}
\]
\wtrap{That equivalence needs the ``one sign change'' proviso. A cash flow
that turns negative again later can have several IRRs or none, and then only
the NPV is meaningful. When the two criteria disagree, rely on the NPV.}

\wstep{Step 8. The three formulas this example used, in one place.}
\[
  \text{annuity PV}=R\,\frac{1-(1+i)^{-n}}{i},
  \qquad
  \text{annuity FV}=R\,\frac{(1+i)^{n}-1}{i},
  \qquad
  \text{perpetuity}=\frac{R}{i} .
\]
All three are geometric series from \chref{foundations}{2}: the perpetuity is the
infinite case, and letting $n\to\infty$ in the first indeed gives $R/i$.
\end{worked}

\section{Practice Problems}

\begin{problems}
  \item A capital of $K_{0}=1000$ is invested for $10$ years at $5\%$
        p.a. Compute the end value under (a) simple interest,
        (b) annual compound interest.
  \item Compute the end value of $K_{0}=2500$ after $5$ years at $4\%$
        p.a.\ compound interest.
  \item What is the present value of a payment of $10\,000$ due in
        $8$ years, at $5\%$ p.a.?
  \item A zero-coupon bond repays $5000$ in $6$ years. At a market rate
        of $2\%$ p.a., what price (present value) is fair today?
  \item A bank offers a nominal rate of $8\%$ p.a.\ compounded
        quarterly. Compute the effective annual rate and decide whether
        a saver should prefer this account or one paying $8.25\%$ with
        annual compounding.
  \item Compute the effective annual rate of a nominal $12\%$ p.a.\
        compounded monthly.
  \item How long does it take to double a capital at $5\%$ p.a.\
        compound interest? (Solve $1.05^{n}=2$ with logarithms.)
  \item How long does it take to double a capital at $2\%$ p.a.\
        compound interest?
  \item $K_{0}=1000$ is invested for $10$ years at $r=5\%$ p.a.\ with
        \emph{continuous} compounding, $K(t)=K_{0}\mathrm{e}^{rt}$.
        Compute the end value and compare with annual compounding from
        Problem~1.
  \item Savings plan: you save $200$ per month; for clean numbers assume
        the whole $R=2400$ is deposited at the end of each year. At
        $5\%$ p.a., what has the plan grown to immediately after the
        $10$th deposit?
  \item Repeat Problem~10 with the deposits made at the
        \emph{beginning} of each year (annuity-due).
  \item An ordinary annuity pays $R=1000$ at the end of each year for
        $5$ years; $i=10\%$. Compute its present value.
  \item Perpetuity: what is the present value of $R=500$ per year
        forever at $i=4\%$?
  \item A perpetual bond paying $1500$ per year costs $25\,000$. Which
        interest rate does this price imply?
  \item A debt of $S=100\,000$ is to be repaid by equal annual payments
        over $20$ years at $5\%$ p.a. Compute the annuity payment.
  \item A loan of $10\,000$ at $10\%$ p.a.\ is repaid by three equal
        annual payments. Compute the payment and set up the full
        amortization schedule (interest, principal and remaining debt
        for each of the three years).
  \item Two projects each cost $1000$ today. Project~A returns $600$ at
        the end of each of years $1$ and $2$; project~B returns a single
        $1250$ at the end of year~$2$. Compare the net present values at
        $i=10\%$: which project should be chosen?
  \item A project costs $2000$ today and returns $800$ at the end of
        each of years $1$--$3$. Compute the NPV at $i=5\%$. Is the
        project worthwhile?
  \item Three deposits of $100$, made at the ends of years $1,2,3$, have
        grown to exactly $331$ immediately after the last deposit. Use
        the geometric sum to set up an equation for $q=1+i$ and solve
        for the interest rate $i$.
  \item $K_{0}=1000$ is invested for $3$ years at $4\%$ p.a.\ and then
        for $2$ further years at $10\%$ p.a.\ (compound interest
        throughout). Compute the end value.
\end{problems}

\section{Solutions}

Money is rounded to two decimal places at the final step (and, in
amortization schedules, to the cent in every row); growth and discount
factors are carried at full precision.

\begin{solutions}
  \item (a) Simple interest: $K_{n}=K_{0}(1+n\,i)
        =1000\,(1+10\cdot 0.05)=1000\cdot 1.5=1500.00$.\\
        (b) Compound interest: $K_{n}=K_{0}(1+i)^{n}
        =1000\cdot 1.05^{10}=1000\cdot 1.6288946=1628.89$.
        Compounding earns $128.89$ more.

  \item $K_{n}=K_{0}(1+i)^{n}=2500\cdot 1.04^{5}$. Since
        $1.04^{5}=1.2166529$,
        $K_{5}=2500\cdot 1.2166529=3041.63$.

  \item $K_{0}=K_{n}(1+i)^{-n}=\dfrac{10\,000}{1.05^{8}}$. Since
        $1.05^{8}=1.4774554$,
        $K_{0}=\dfrac{10\,000}{1.4774554}=6768.39$.

  \item $K_{0}=\dfrac{5000}{1.02^{6}}$. Since $1.02^{6}=1.1261624$,
        $K_{0}=\dfrac{5000}{1.1261624}=4439.86$.

  \item $i_{\mathrm{eff}}=\left(1+\frac{j}{m}\right)^{m}-1
        =\left(1+\frac{0.08}{4}\right)^{4}-1=1.02^{4}-1$. Since
        $1.02^{4}=1.08243216$,
        $i_{\mathrm{eff}}=0.08243216\approx 8.2432\%$. As
        $8.2432\%<8.25\%$, the account with annual compounding is
        (slightly) better for the saver.

  \item $i_{\mathrm{eff}}=\left(1+\frac{0.12}{12}\right)^{12}-1
        =1.01^{12}-1=1.1268250-1=0.1268250\approx 12.6825\%$.

  \item $1.05^{n}=2\Rightarrow n\ln 1.05=\ln 2\Rightarrow
        n=\dfrac{\ln 2}{\ln 1.05}=\dfrac{0.6931472}{0.0487902}
        \approx 14.21$ years. Since interest is credited only at year
        ends, the capital has actually doubled after $15$ full years
        ($1.05^{14}=1.9799<2<2.0789=1.05^{15}$).

  \item $n=\dfrac{\ln 2}{\ln 1.02}=\dfrac{0.6931472}{0.0198026}
        \approx 35.00$ years. (Because $n=35.0028>35$, with annual
        crediting the capital first exceeds the double after $36$ full
        years: $1.02^{35}\approx 1.9999<2$.)

  \item $K(10)=K_{0}\mathrm{e}^{rt}=1000\,\mathrm{e}^{0.05\cdot 10}
        =1000\,\mathrm{e}^{0.5}=1000\cdot 1.6487213=1648.72$.
        Continuous compounding beats the annual value $1628.89$ from
        Problem~1 by $19.83$.

  \item Monthly saving $200$ means $R=12\cdot 200=2400$ per year. End
        value of the ordinary annuity:
        \[
          E_{10}=R\,\frac{q^{10}-1}{i}
                =2400\cdot\frac{1.05^{10}-1}{0.05}
                =2400\cdot\frac{0.6288946}{0.05}
                =2400\cdot 12.5778925=30\,186.94 .
        \]

  \item Annuity-due: every deposit earns one extra year, so
        $E_{10}^{\mathrm{due}}=q\,E_{10}
        =1.05\cdot 30\,186.94=31\,696.29$.

  \item $B_{n}=R\,\dfrac{1-q^{-n}}{i}
        =1000\cdot\dfrac{1-1.1^{-5}}{0.1}$. Since $1.1^{5}=1.61051$,
        $1.1^{-5}=0.6209213$ and
        \[
          B_{5}=1000\cdot\frac{1-0.6209213}{0.1}
               =1000\cdot 3.7907868=3790.79 .
        \]

  \item $B_{\infty}=\dfrac{R}{i}=\dfrac{500}{0.04}=12\,500.00$.

  \item From $B_{\infty}=R/i$: $i=\dfrac{R}{B_{\infty}}
        =\dfrac{1500}{25\,000}=0.06=6\%$.

  \item $A=S\,\dfrac{i\,q^{n}}{q^{n}-1}
        =100\,000\cdot\dfrac{0.05\cdot 1.05^{20}}{1.05^{20}-1}$. Since
        $1.05^{20}=2.6532977$,
        \[
          A=100\,000\cdot\frac{0.05\cdot 2.6532977}{1.6532977}
           =100\,000\cdot\frac{0.1326649}{1.6532977}
           =100\,000\cdot 0.0802426=8024.26 .
        \]

  \item Payment: $A=S\,\dfrac{i\,q^{n}}{q^{n}-1}
        =10\,000\cdot\dfrac{0.1\cdot 1.331}{0.331}
        =\dfrac{1331}{0.331}=4021.15$. Schedule (interest $=10\%$ of the
        debt at start, principal $=A-{}$interest, all rounded to the
        cent):

        \begin{center}
        \begin{tabular}{c r r r r}
        \hline
        Year & Debt at start & Interest & Principal & Debt at end\\
        \hline
        $1$ & $10\,000.00$ & $1000.00$ & $3021.15$ & $6978.85$\\
        $2$ & $6978.85$ & $697.89$ & $3323.26$ & $3655.59$\\
        $3$ & $3655.59$ & $365.56$ & $3655.59$ & $0.00$\\
        \hline
        \end{tabular}
        \end{center}

        The interest share falls ($1000.00$, $697.89$, $365.56$) while
        the principal share grows, and the last payment
        $365.56+3655.59=4021.15$ closes the loan exactly.

  \item $\mathrm{NPV}=\sum_{t}C_{t}q^{-t}$ with $q=1.1$:
        \[
          \mathrm{NPV}_{A}=-1000+\frac{600}{1.1}+\frac{600}{1.1^{2}}
          =-1000+545.45+495.87=41.32 ,
        \]
        \[
          \mathrm{NPV}_{B}=-1000+\frac{1250}{1.1^{2}}
          =-1000+\frac{1250}{1.21}=-1000+1033.06=33.06 .
        \]
        Both are worthwhile, but $\mathrm{NPV}_{A}>\mathrm{NPV}_{B}$:
        choose project~A, even though B's undiscounted total
        ($1250$) exceeds A's ($1200$), A pays earlier.

  \item The three equal returns form an ordinary annuity, so
        \[
          \mathrm{NPV}=-2000+800\cdot\frac{1-1.05^{-3}}{0.05}
          =-2000+800\cdot\frac{1-0.8638376}{0.05}
          =-2000+800\cdot 2.7232480 ,
        \]
        using $1.05^{3}=1.157625$. Hence
        $\mathrm{NPV}=-2000+2178.60=178.60>0$: the project is
        worthwhile.

  \item End value of the three deposits (geometric sum):
        $100\,(q^{2}+q+1)=331$, i.e.\ $q^{2}+q+1=3.31$, so
        $q^{2}+q-2.31=0$. The quadratic formula gives
        \[
          q=\frac{-1+\sqrt{1+4\cdot 2.31}}{2}
           =\frac{-1+\sqrt{10.24}}{2}=\frac{-1+3.2}{2}=1.1
        \]
        (the negative root is meaningless). Hence $i=q-1=0.1=10\%$.
        Check: $100\,(1.21+1.1+1)=331$.

  \item Growth factors multiply:
        $K_{5}=K_{0}\cdot 1.04^{3}\cdot 1.1^{2}
        =1000\cdot 1.124864\cdot 1.21=1361.09$.
\end{solutions}

\section{Additional Problems}

\begin{problems}
  \item Compute the end value of $K_{0}=8000$ after $4$ years at $5\%$
        p.a.\ compound interest.
  \item $K_{0}=5000$ is invested for $10$ years at $6\%$ p.a. Compute
        the end value under simple and under compound interest, and the
        difference.
  \item What is the present value of $20\,000$ due in $10$ years at
        $10\%$ p.a.?
  \item A nominal rate of $4\%$ p.a.\ is compounded quarterly. Compute
        the effective annual rate and compare with an account paying
        $4.05\%$ annually: which should a saver choose?
  \item How long does it take to double a capital (a) at $10\%$ p.a.\
        with annual compounding, (b) at $r=10\%$ with continuous
        compounding $K(t)=K_{0}\mathrm{e}^{rt}$?
  \item A savings plan deposits $R=1000$ at the end of each year for
        $5$ years at $4\%$ p.a. Compute the end value immediately after
        the last deposit.
  \item A pension pays $6000$ at the end of each year for $10$ years;
        $i=5\%$. Compute its present value.
  \item An endowment is to fund a prize of $3000$ per year forever at
        $i=6\%$. How much capital is required today?
  \item A loan of $50\,000$ at $4\%$ p.a.\ is repaid by five equal
        annual payments. Compute the payment.
  \item A project costs $5000$ today and returns $2000$ at the end of
        each of years $1$--$3$. Compute the NPV at $i=10\%$ and decide.
        Compare with the undiscounted total of the returns.
\end{problems}

\section{Additional Solutions}

\begin{solutions}
  \item $K_{n}=K_{0}(1+i)^{n}=8000\cdot 1.05^{4}
        =8000\cdot 1.21550625=9724.05$.

  \item Simple: $K_{n}=K_{0}(1+n\,i)=5000\,(1+10\cdot 0.06)
        =5000\cdot 1.6=8000.00$. Compound:
        $K_{n}=5000\cdot 1.06^{10}=5000\cdot 1.7908477=8954.24$.
        Compounding earns $8954.24-8000.00=954.24$ more.

  \item $K_{0}=\dfrac{20\,000}{1.1^{10}}=\dfrac{20\,000}{2.5937425}
        =7710.87$.

  \item $i_{\mathrm{eff}}=\left(1+\frac{0.04}{4}\right)^{4}-1
        =1.01^{4}-1=1.04060401-1=0.04060401\approx 4.0604\%$. As
        $4.0604\%>4.05\%$, the quarterly account is (slightly) better
        for the saver.

  \item (a) $1.1^{n}=2\Rightarrow n=\dfrac{\ln 2}{\ln 1.1}
        =\dfrac{0.6931472}{0.0953102}\approx 7.27$ years.\\
        (b) $\mathrm{e}^{0.1t}=2\Rightarrow t=\dfrac{\ln 2}{0.1}
        =6.93$ years, continuous compounding doubles faster.

  \item $E_{5}=R\,\dfrac{q^{5}-1}{i}
        =1000\cdot\dfrac{1.04^{5}-1}{0.04}
        =1000\cdot\dfrac{0.2166529}{0.04}
        =1000\cdot 5.4163226=5416.32$.

  \item $B_{10}=R\,\dfrac{1-q^{-10}}{i}
        =6000\cdot\dfrac{1-1.05^{-10}}{0.05}$. Since
        $1.05^{10}=1.6288946$, $1.05^{-10}=0.6139133$ and
        \[
          B_{10}=6000\cdot\frac{1-0.6139133}{0.05}
                =6000\cdot 7.7217349=46\,330.41 .
        \]

  \item $B_{\infty}=\dfrac{R}{i}=\dfrac{3000}{0.06}=50\,000.00$.

  \item $A=S\,\dfrac{i}{1-q^{-n}}
        =50\,000\cdot\dfrac{0.04}{1-1.04^{-5}}$. Since
        $1.04^{5}=1.2166529$, $1.04^{-5}=0.8219271$ and
        \[
          A=50\,000\cdot\frac{0.04}{0.1780729}
           =50\,000\cdot 0.2246271=11\,231.36 .
        \]

  \item $\mathrm{NPV}=-5000+2000\cdot\dfrac{1-1.1^{-3}}{0.1}$. Since
        $1.1^{3}=1.331$, $1.1^{-3}=0.7513148$ and
        \[
          \mathrm{NPV}=-5000+2000\cdot\frac{1-0.7513148}{0.1}
                      =-5000+2000\cdot 2.4868520
        \]
        \[
          \phantom{\mathrm{NPV}}=-5000+4973.70=-26.30<0 :
        \]
        reject the project. The undiscounted total $3\cdot 2000=6000$
        exceeds the outlay $5000$, but discounting at $10\%$ turns the
        surplus into a loss, totals without discounting mislead.
\end{solutions}

\chapter{Limits of Functions and Continuity}

The limits of sequences studied so far let $n$ run through the natural
numbers. This closing chapter transfers the idea to functions of a
continuous variable: what value does $f(x)$ approach as $x$ approaches a
point $a$, or as $x$ grows beyond all bounds? The answer leads directly
to \emph{continuity}, the property that makes a function
graph drawable without lifting the pen, and to the Intermediate Value
Theorem, the tool that guarantees roots and market equilibria exist and
tells us how to hunt them down numerically.

The treatment stays computational: limits are evaluated with the same
limit rules as for sequences, and the indeterminate form $0/0$ is resolved
by algebra (factoring, expanding, or multiplying by a conjugate)
never by differentiation, which only enters in \volref{derivatives}. Economic
examples run through the chapter: price schedules with quantity discounts
are natural piecewise functions, and their discontinuities have a concrete
monetary meaning.

\section{Core Rules}

\begin{rulebox}{Limit of a function; one-sided limits}
$\displaystyle\lim_{x\to a}f(x)=L$ means: as $x$ gets
arbitrarily close to $a$ (from either side, without ever equalling $a$) the values $f(x)$ get arbitrarily close to $L$. The point $a$ itself
need not lie in the domain of $f$, and $f(a)$, if it exists, plays no role.

The \textbf{one-sided limits}
\[
  \lim_{x\to a^{-}}f(x)
  \qquad\text{and}\qquad
  \lim_{x\to a^{+}}f(x)
\]
approach $a$ from the left ($x<a$) and from the right ($x>a$) only.
\textbf{The two-sided limit exists exactly when both one-sided limits
exist and are equal}, and then all three coincide. For piecewise-defined
functions, always compute the two one-sided limits separately.
\end{rulebox}

\begin{rulebox}{Limit rules}
If $\lim_{x\to a}f(x)=L$ and $\lim_{x\to a}g(x)=M$ (both finite), then
\begin{itemize}[leftmargin=1.6em,itemsep=0.3ex]
  \item $\lim\limits_{x\to a}\bigl(f(x)\pm g(x)\bigr)=L\pm M$,
  \item $\lim\limits_{x\to a}\bigl(f(x)\cdot g(x)\bigr)=L\cdot M$,
        \quad in particular $\lim\limits_{x\to a}c\,f(x)=cL$,
  \item $\lim\limits_{x\to a}\dfrac{f(x)}{g(x)}=\dfrac{L}{M}$ \quad
        provided $M\neq 0$,
  \item $\lim\limits_{x\to a}f(x)^{k}=L^{k}$ and
        $\lim\limits_{x\to a}\sqrt[k]{f(x)}=\sqrt[k]{L}$ (where defined).
\end{itemize}
These are the same rules as for sequences, with $x\to a$ in place of
$n\to\infty$; they also hold for $x\to\pm\infty$ and for one-sided limits.
For all the standard functions below, limits at points of the domain are
found by \textbf{direct substitution}: $\lim_{x\to a}f(x)=f(a)$.
\end{rulebox}

\begin{rulebox}{The indeterminate form $0/0$}
If numerator and denominator both tend to $0$, the quotient rule does not
apply: the limit can still be anything. Two algebraic escapes:
\begin{itemize}[leftmargin=1.6em,itemsep=0.3ex]
  \item \textbf{Factor and cancel:} if a polynomial vanishes at $x=a$,
        then $(x-a)$ is one of its factors. Factor numerator and
        denominator, cancel the common factor $(x-a)$, then substitute.
        Example:
        \[
          \lim_{x\to 2}\frac{x^{2}-4}{x-2}
          =\lim_{x\to 2}\frac{(x-2)(x+2)}{x-2}
          =\lim_{x\to 2}\,(x+2)=4 .
        \]
  \item \textbf{Conjugate trick for roots:} expand the fraction with the
        conjugate $\sqrt{u}+v$ of a root expression $\sqrt{u}-v$ and use
        $(\sqrt{u}-v)(\sqrt{u}+v)=u-v^{2}$. Example:
        \[
          \lim_{x\to 0}\frac{\sqrt{x+9}-3}{x}
          =\lim_{x\to 0}\frac{x}{x\bigl(\sqrt{x+9}+3\bigr)}
          =\lim_{x\to 0}\frac{1}{\sqrt{x+9}+3}=\frac{1}{6} .
        \]
\end{itemize}
\end{rulebox}

\begin{rulebox}{Limits at $\pm\infty$; horizontal and vertical asymptotes}
For a rational function $\dfrac{p(x)}{q(x)}$ compare the degrees: divide
numerator and denominator by the highest power of $x$ in the denominator.
\begin{itemize}[leftmargin=1.6em,itemsep=0.3ex]
  \item $\deg p<\deg q$: \quad limit $0$ as $x\to\pm\infty$,
  \item $\deg p=\deg q$: \quad limit $=$ ratio of the leading
        coefficients,
  \item $\deg p>\deg q$: \quad the limit is $+\infty$ or $-\infty$
        (sign from the leading terms).
\end{itemize}
The line $y=L$ is a \textbf{horizontal asymptote} if $f(x)\to L$ as $x\to\infty$ or $x\to-\infty$. The line
$x=a$ is a \textbf{vertical asymptote} if at
least one one-sided limit at $a$ is $+\infty$ or $-\infty$; for rational
functions this happens where the denominator vanishes but the numerator
does not \emph{after cancelling common factors}. A common factor that
cancels completely produces a hole in the graph, not an asymptote.
\end{rulebox}

\begin{rulebox}{Sandwich rule}
If $g(x)\le f(x)\le h(x)$ for all $x$ near $a$ and
\[
  \lim_{x\to a}g(x)=\lim_{x\to a}h(x)=L,
  \qquad\text{then}\qquad
  \lim_{x\to a}f(x)=L .
\]
The same holds for $x\to\pm\infty$. The standard application is
``bounded times vanishing'': since $-1\le\sin t\le 1$ for all $t$,
\[
  -x^{2}\;\le\; x^{2}\sin\frac{1}{x}\;\le\; x^{2}
  \quad\Longrightarrow\quad
  \lim_{x\to 0}x^{2}\sin\frac{1}{x}=0 .
\]
\end{rulebox}

\begin{rulebox}{Continuity}
$f$ is \textbf{continuous at $a$} if three things hold:
$f(a)$ is defined, $\lim_{x\to a}f(x)$ exists, and
\[
  \lim_{x\to a}f(x)=f(a) .
\]
$f$ is \textbf{continuous on an interval} if it is continuous at every
point of the interval. All the standard functions are continuous wherever
they are defined: polynomials and rational functions, roots $\sqrt[k]{x}$,
powers $x^{r}$, the absolute value $|x|$, exponentials $\mathrm{e}^{x}$
and $a^{x}$, the logarithm $\ln x$, and $\sin x$, $\cos x$. Moreover,
\textbf{sums, differences, products, quotients (where the denominator is
nonzero), and compositions of continuous functions are continuous}. In
practice, discontinuities can therefore only occur at points where a
formula breaks down (denominator $=0$, log of $0$, \dots) or where a
piecewise definition switches branches.
\end{rulebox}

\begin{rulebox}{Classification of discontinuities}
Suppose $f$ is not continuous at $a$. Compute the one-sided limits.
\begin{itemize}[leftmargin=1.6em,itemsep=0.3ex]
  \item \textbf{Removable}: the
        two-sided limit $L$ exists (finite), but $f(a)\neq L$ or $f(a)$
        is undefined. Redefining $f(a):=L$ repairs the function:
        the graph has a single hole. Typical source: a cancelled $0/0$.
  \item \textbf{Jump}: both one-sided limits exist
        and are finite but differ. The jump has size
        $\lim_{x\to a^{+}}f(x)-\lim_{x\to a^{-}}f(x)$.
  \item \textbf{Pole}: at least one one-sided limit is
        $+\infty$ or $-\infty$; the line $x=a$ is a vertical asymptote.
\end{itemize}
Piecewise price schedules are the standard economic example: if a
quantity discount applies to \emph{all} units once the order passes a
threshold, the total bill jumps \emph{downward} at the threshold: a
jump discontinuity that makes some order sizes irrational to choose. If
the discount applies only to the units \emph{above} the threshold, the
bill is continuous.
\end{rulebox}

\begin{rulebox}{Intermediate Value Theorem (IVT)}
Let $f$ be continuous on $[a,b]$. Then $f$
attains every value between $f(a)$ and $f(b)$. In particular:
\[
  f(a)\cdot f(b)<0
  \quad\Longrightarrow\quad
  f(x^{*})=0 \text{ for some } x^{*}\in(a,b).
\]
This is the existence tool for roots and for market equilibria (apply it
to excess demand $D(p)-S(p)$). It also drives the \textbf{bisection
method}: evaluate $f$ at the midpoint $m=\tfrac{a+b}{2}$ and keep the
half-interval on which $f$ still changes sign. Each step halves the
bracket, so after $n$ steps the root is located to within
$(b-a)/2^{n}$. Continuity is essential: a function with a jump can change
sign without ever being zero.
\end{rulebox}

\begin{rulebox}{Weierstrass: extreme values on $[a,b]$}
A function continuous on a \emph{closed, bounded} interval $[a,b]$
attains a maximum and a minimum there.
This is the one-variable case of the extreme value theorem that later
volumes drill in several variables; here it guarantees, before any
optimisation technique is available, that ``best'' and ``worst'' values
exist at all. Both hypotheses matter: $f(x)=1/x$ on $(0,1]$ has no
maximum, and $f(x)=x$ on $[0,\infty)$ has none either.
\end{rulebox}

\section{Worked Example}

\begin{worked}{Evaluate $\lim_{x\to0}\dfrac{\sqrt{x+4}-2}{x}$, choose $a$ and $b$ making a piecewise function continuous, and classify the discontinuities of three examples}
\wstep{Step 1. Substitute first --- and read what the substitution says.}
At $x=0$ the expression is $\frac{2-2}{0}=\frac00$. That is \textbf{not} an
answer and \textbf{not} a proof that the limit fails to exist; it is a
signal that numerator and denominator share a vanishing factor, which
algebra must expose.

\wstep{Step 2. Rationalise the numerator.}
A difference of square roots is cleared by multiplying by its conjugate:
\[
  \frac{\sqrt{x+4}-2}{x}\cdot\frac{\sqrt{x+4}+2}{\sqrt{x+4}+2}
  =\frac{(x+4)-4}{x\bigl(\sqrt{x+4}+2\bigr)}
  =\frac{x}{x\bigl(\sqrt{x+4}+2\bigr)} .
\]
The factor $x$ has been made visible --- which was the entire purpose.
\wtrap{Multiply the \emph{numerator} by its conjugate when the roots are on
top. The reflex of rationalising the denominator is unhelpful here; the
goal is not tidiness but exposing the common factor.}

\wstep{Step 3. Cancel, then substitute.}
For $x\neq0$ --- and $x\neq0$ is exactly the situation a limit concerns ---
\[
  \frac{x}{x(\sqrt{x+4}+2)}=\frac{1}{\sqrt{x+4}+2}
  \xrightarrow[x\to0]{}\frac{1}{2+2}=\boxed{\tfrac14}.
\]
Numerical check: at $x=0.01$ the original quotient is $0.2499\ldots$
\checkmark

\wstep{Step 4. Now make $f$ continuous at $x=0$.}
Define
\[
  f(x)=\begin{cases}\dfrac{\sqrt{x+4}-2}{x}, & x\neq0,\\[1.2ex] a, & x=0.\end{cases}
\]
Continuity at a point requires three things: $f(0)$ exists, the limit
exists, and they are \emph{equal}. The first two hold, so the condition is
$a=\tfrac14$. This is a \textbf{removable} discontinuity: one value, chosen
correctly, repairs it.

\wstep{Step 5. A two-sided version, where both a limit and a slope must match.}
\[
  g(x)=\begin{cases}x^{2}, & x\le1,\\ ax+b, & x>1.\end{cases}
\]
Continuity at $1$ needs the one-sided limits to agree:
\[
  \lim_{x\to1^{-}}g=1,\qquad \lim_{x\to1^{+}}g=a+b
  \quad\Longrightarrow\quad a+b=1 .
\]
If differentiability is also demanded, the slopes must match too:
$2x|_{x=1}=2$ from the left and $a$ from the right, so $a=2$ and hence
$b=-1$.
\wtrap{Continuity gives \emph{one} equation, differentiability a second.
With two unknowns you need both conditions --- and note the order: a
function can be continuous without being differentiable ($|x|$ at $0$), but
never differentiable without being continuous.}

\wstep{Step 6. Classify three discontinuities.}
\begin{center}\footnotesize
\begin{tabular}{lll}
  \hline
  function & at & type\\
  \hline
  $\dfrac{x^{2}-1}{x-1}$ & $x=1$ & removable: both one-sided limits are $2$\\[0.8ex]
  $\dfrac{1}{x-2}$ & $x=2$ & infinite (pole): $\to-\infty$ left, $+\infty$ right\\[0.8ex]
  $\sgn(x)$ & $x=0$ & jump: limits $-1$ and $+1$, both finite, unequal\\
  \hline
\end{tabular}
\end{center}
\wtrap{All three have ``a problem at a point'', and they behave completely
differently. Decide by computing the two one-sided limits: equal and finite
$\Rightarrow$ removable; finite but different $\Rightarrow$ jump; at least
one infinite $\Rightarrow$ pole. Only the removable kind can be patched.}

\wstep{Step 7. Continuity earns two theorems --- use them.}
\emph{Intermediate Value Theorem.} $h(x)=x^{3}+x-1$ is continuous, with
$h(0)=-1<0$ and $h(1)=1>0$. So $h$ takes the value $0$ somewhere in
$(0,1)$: the equation has a root, proved without solving it. This is how
every existence claim about an equilibrium price is made.

\emph{Weierstrass.} A continuous function on a \textbf{closed and bounded}
interval attains a maximum and a minimum. Both hypotheses are load-bearing:
$1/x$ on $(0,1]$ is continuous but unbounded, because the interval is not
closed; $x$ on $[0,\infty)$ has no maximum, because the interval is not
bounded.
\wtrap{In \volref{optimization} the existence half of every optimisation problem is
Weierstrass, and the usual failure is a feasible set that is open ($x>0$
instead of $x\ge0$) or unbounded. When an optimisation problem ``has no
solution'', check those two words before checking the calculus.}
\end{worked}

\section{Practice Problems}

Evaluate each limit or answer the question posed. Where a limit fails to
exist, say so and give the one-sided limits; write $+\infty$ or $-\infty$
where appropriate.

\begin{problems}
  \item $\displaystyle\lim_{x\to 2}\frac{x^{2}+1}{x+3}$
  \item $\displaystyle\lim_{x\to 2}\frac{x^{2}-4}{x-2}$
  \item $\displaystyle\lim_{x\to -3}\frac{x^{2}+5x+6}{x+3}$
  \item $\displaystyle\lim_{x\to 1}\frac{x^{3}-1}{x-1}$
  \item $\displaystyle\lim_{x\to 2}\frac{x^{2}-5x+6}{x^{2}-4}$
  \item $\displaystyle\lim_{x\to 4}\frac{\sqrt{x}-2}{x-4}$
  \item $\displaystyle\lim_{x\to 0}\frac{\sqrt{x+9}-3}{x}$
  \item For $f(x)=\dfrac{|x-2|}{x-2}$ compute
        $\displaystyle\lim_{x\to 2^{-}}f(x)$ and
        $\displaystyle\lim_{x\to 2^{+}}f(x)$. Does
        $\displaystyle\lim_{x\to 2}f(x)$ exist?
  \item For
        \[
          f(x)=
          \begin{cases}
            x^{2}, & x<1,\\
            3-x,   & x\ge 1,
          \end{cases}
        \]
        compute both one-sided limits at $x=1$. Does the two-sided limit
        exist? Classify the discontinuity, if any.
  \item $\displaystyle\lim_{x\to\infty}\frac{3x^{2}-5x+1}{2x^{2}+7}$
        \quad and \quad
        $\displaystyle\lim_{x\to-\infty}\frac{3x^{2}-5x+1}{2x^{2}+7}$
  \item $\displaystyle\lim_{x\to\infty}\frac{2x+1}{x^{2}-4}$. Which
        horizontal asymptote does this give?
  \item $\displaystyle\lim_{x\to\infty}\frac{x^{3}-2x}{5x^{2}+1}$
        \quad and \quad
        $\displaystyle\lim_{x\to-\infty}\frac{x^{3}-2x}{5x^{2}+1}$
  \item Find all horizontal and vertical asymptotes of
        $f(x)=\dfrac{2x^{2}-8}{x^{2}-1}$, including the one-sided
        behaviour at each vertical asymptote.
  \item Find all asymptotes of $f(x)=\dfrac{x^{2}-9}{x-3}$ and classify
        the discontinuity at $x=3$.
  \item Use the sandwich rule to compute
        $\displaystyle\lim_{x\to 0}\,x^{2}\sin\frac{1}{x}$.
  \item Find the constant $c$ that makes
        \[
          f(x)=
          \begin{cases}
            cx+3,    & x<2,\\
            x^{2}+1, & x\ge 2,
          \end{cases}
        \]
        continuous at $x=2$ (and hence on all of $\R$).
  \item Classify all discontinuities of
        $f(x)=\dfrac{x+2}{x^{2}-4}$.
  \item A supplier offers a quantity discount on the \emph{whole} order:
        \[
          C(q)=
          \begin{cases}
            12q, & 0\le q\le 50,\\
            9q,  & q>50,
          \end{cases}
        \]
        where $C(q)$ is the total bill for $q$ units.
        (a)~Compute the one-sided limits of $C$ at $q=50$ and classify
        the discontinuity.
        (b)~The supplier switches to $C(q)=9q+F$ for $q>50$ with a fixed
        fee $F$. Which $F$ makes $C$ continuous at $q=50$, and what does
        the continuous schedule mean economically?
  \item Show with the Intermediate Value Theorem that
        $f(x)=x^{3}-3x+1$ has a root in $[0,1]$, then perform one
        bisection step: evaluate $f$ at the midpoint and state the new,
        halved bracket.
  \item On which set is
        $g(x)=\dfrac{\sqrt{x+2}}{x-1}$
        continuous? Justify using the continuity of standard functions
        and of compositions and quotients.
\end{problems}

\section{Solutions}

\begin{solutions}
  \item The denominator does not vanish at $x=2$, so substitute
        directly:
        \[
          \lim_{x\to 2}\frac{x^{2}+1}{x+3}=\frac{4+1}{2+3}=1 .
        \]

  \item At $x=2$ both numerator and denominator are $0$: the form is
        $0/0$. Factor and cancel:
        \[
          \frac{x^{2}-4}{x-2}=\frac{(x-2)(x+2)}{x-2}=x+2
          \quad (x\neq 2),
          \qquad
          \lim_{x\to 2}\,(x+2)=4 .
        \]

  \item Form $0/0$. Since $x^{2}+5x+6=(x+2)(x+3)$,
        \[
          \lim_{x\to -3}\frac{(x+2)(x+3)}{x+3}
          =\lim_{x\to -3}\,(x+2)=-1 .
        \]

  \item Form $0/0$. Use $x^{3}-1=(x-1)(x^{2}+x+1)$:
        \[
          \lim_{x\to 1}\frac{(x-1)(x^{2}+x+1)}{x-1}
          =\lim_{x\to 1}\,(x^{2}+x+1)=1+1+1=3 .
        \]

  \item Form $0/0$; both quadratics contain the factor $(x-2)$:
        \[
          \frac{x^{2}-5x+6}{x^{2}-4}
          =\frac{(x-2)(x-3)}{(x-2)(x+2)}
          =\frac{x-3}{x+2}
          \quad (x\neq 2),
        \]
        so the limit is $\dfrac{2-3}{2+2}=-\dfrac{1}{4}$.

  \item Form $0/0$. Expand with the conjugate $\sqrt{x}+2$:
        \[
          \frac{\sqrt{x}-2}{x-4}
          =\frac{(\sqrt{x}-2)(\sqrt{x}+2)}{(x-4)(\sqrt{x}+2)}
          =\frac{x-4}{(x-4)(\sqrt{x}+2)}
          =\frac{1}{\sqrt{x}+2},
        \]
        hence the limit is $\dfrac{1}{\sqrt{4}+2}=\dfrac{1}{4}$.

  \item Form $0/0$. Conjugate $\sqrt{x+9}+3$:
        \[
          \frac{\sqrt{x+9}-3}{x}
          =\frac{(x+9)-9}{x\bigl(\sqrt{x+9}+3\bigr)}
          =\frac{1}{\sqrt{x+9}+3}
          \;\xrightarrow[x\to 0]{}\;
          \frac{1}{3+3}=\frac{1}{6} .
        \]

  \item For $x<2$ we have $|x-2|=-(x-2)$, so $f(x)=-1$; for $x>2$ we
        have $|x-2|=x-2$, so $f(x)=+1$. Hence
        \[
          \lim_{x\to 2^{-}}f(x)=-1,
          \qquad
          \lim_{x\to 2^{+}}f(x)=+1 .
        \]
        The one-sided limits differ, so $\lim_{x\to 2}f(x)$ does not
        exist ($f$ has a jump of size $2$ at $x=2$).

  \item Left branch: $\lim_{x\to 1^{-}}x^{2}=1$. Right branch:
        $\lim_{x\to 1^{+}}(3-x)=2$, and $f(1)=3-1=2$. The one-sided
        limits are finite but different, so the two-sided limit does not
        exist: $f$ has a \textbf{jump discontinuity} at $x=1$ with jump
        $2-1=1$.

  \item Degrees are equal, so the limit is the ratio of the leading
        coefficients. Dividing by $x^{2}$:
        \[
          \frac{3x^{2}-5x+1}{2x^{2}+7}
          =\frac{3-\frac{5}{x}+\frac{1}{x^{2}}}{2+\frac{7}{x^{2}}}
          \;\longrightarrow\;\frac{3}{2}
          \qquad (x\to\infty\ \text{and}\ x\to-\infty).
        \]
        Both limits equal $\tfrac{3}{2}$.

  \item The numerator has smaller degree than the denominator, so
        \[
          \lim_{x\to\infty}\frac{2x+1}{x^{2}-4}
          =\lim_{x\to\infty}
           \frac{\frac{2}{x}+\frac{1}{x^{2}}}{1-\frac{4}{x^{2}}}
          =\frac{0}{1}=0 .
        \]
        The horizontal asymptote is $y=0$ (the $x$-axis).

  \item The numerator has larger degree. Divide by $x^{2}$:
        \[
          \frac{x^{3}-2x}{5x^{2}+1}
          =\frac{x-\frac{2}{x}}{5+\frac{1}{x^{2}}} .
        \]
        The denominator tends to $5$ while the numerator tends to
        $+\infty$ resp.\ $-\infty$, so
        \[
          \lim_{x\to\infty}\frac{x^{3}-2x}{5x^{2}+1}=+\infty,
          \qquad
          \lim_{x\to-\infty}\frac{x^{3}-2x}{5x^{2}+1}=-\infty .
        \]

  \item \textbf{Horizontal:} equal degrees, ratio of leading
        coefficients $\tfrac{2}{1}=2$, so $y=2$ as $x\to\pm\infty$.
        \textbf{Vertical:} the denominator $x^{2}-1=(x-1)(x+1)$
        vanishes at $x=\pm 1$; the numerator $2x^{2}-8$ equals $-6\neq 0$
        at both points, so no cancellation: both are poles. Near $x=1$
        the numerator is close to $-6$, and $(x-1)(x+1)$ is a small
        positive number for $x\to 1^{+}$, a small negative number for
        $x\to 1^{-}$:
        \[
          \lim_{x\to 1^{+}}f(x)=-\infty,\qquad
          \lim_{x\to 1^{-}}f(x)=+\infty .
        \]
        Near $x=-1$ the numerator is again close to $-6$, while
        $(x-1)(x+1)$ is negative for $x\to-1^{+}$ and positive for
        $x\to-1^{-}$:
        \[
          \lim_{x\to -1^{+}}f(x)=+\infty,\qquad
          \lim_{x\to -1^{-}}f(x)=-\infty .
        \]
        Asymptotes: $y=2$, $x=1$, $x=-1$.

  \item Factor: $x^{2}-9=(x-3)(x+3)$, so for $x\neq 3$
        \[
          f(x)=\frac{(x-3)(x+3)}{x-3}=x+3,
          \qquad
          \lim_{x\to 3}f(x)=6 .
        \]
        The limit exists but $f(3)$ is undefined: the discontinuity at
        $x=3$ is \textbf{removable} (a hole at $(3,6)$), \emph{not} a
        vertical asymptote, the common factor cancels completely.
        Since $f(x)=x+3\to\pm\infty$ as $x\to\pm\infty$, there is no
        horizontal asymptote either: $f$ has no asymptotes at all.

  \item Since $-1\le\sin\frac{1}{x}\le 1$ for all $x\neq 0$ and
        $x^{2}\ge 0$,
        \[
          -x^{2}\;\le\; x^{2}\sin\frac{1}{x}\;\le\; x^{2}.
        \]
        Both bounds tend to $0$ as $x\to 0$, so by the sandwich rule
        \[
          \lim_{x\to 0}\,x^{2}\sin\frac{1}{x}=0 .
        \]

  \item Both branches are polynomials, hence continuous away from
        $x=2$; only the seam matters. One-sided limits at $x=2$:
        \[
          \lim_{x\to 2^{-}}(cx+3)=2c+3,
          \qquad
          \lim_{x\to 2^{+}}(x^{2}+1)=5=f(2).
        \]
        Continuity requires $2c+3=5$, i.e.\ $c=1$. With $c=1$
        both one-sided limits equal $f(2)=5$, so $f$ is continuous on
        $\R$.

  \item Factor the denominator: $x^{2}-4=(x-2)(x+2)$, so for
        $x\neq -2$
        \[
          f(x)=\frac{x+2}{(x-2)(x+2)}=\frac{1}{x-2}.
        \]
        \textbf{At $x=-2$:} the factor cancels and
        $\lim_{x\to-2}f(x)=\frac{1}{-4}=-\frac14$ exists, while $f(-2)$
        is undefined: a \textbf{removable} discontinuity.
        \textbf{At $x=2$:} after cancelling, the denominator still
        vanishes while the numerator is $1$:
        \[
          \lim_{x\to 2^{-}}\frac{1}{x-2}=-\infty,
          \qquad
          \lim_{x\to 2^{+}}\frac{1}{x-2}=+\infty,
        \]
        a \textbf{pole} (vertical asymptote $x=2$).

  \item (a) One-sided limits at the threshold:
        \[
          \lim_{q\to 50^{-}}C(q)=12\cdot 50=600=C(50),
          \qquad
          \lim_{q\to 50^{+}}C(q)=9\cdot 50=450 .
        \]
        Both are finite but different: a \textbf{jump discontinuity}
        with jump $450-600=-150$. The bill \emph{drops} by $150$ as the
        order passes $50$ units: ordering $51$ units costs
        $9\cdot 51=459$, less than the $600$ charged for $50$ units,
        indeed less than the $12\cdot 39=468$ charged for $39$ units.
        No rational customer orders between about $39$ and $50$ units.

        (b) Continuity at $q=50$ requires
        \[
          \lim_{q\to 50^{+}}(9q+F)=450+F\overset{!}{=}600
          \quad\Longrightarrow\quad
          F=150 .
        \]
        Then for $q>50$,
        \[
          C(q)=9q+150=12\cdot 50+9\,(q-50):
        \]
        the first $50$ units cost $12$ each and only the units
        \emph{beyond} the threshold get the discounted price $9$. The
        continuous schedule is exactly the marginal-discount tariff.

  \item $f(x)=x^{3}-3x+1$ is a polynomial, hence continuous on $[0,1]$.
        Endpoint values:
        \[
          f(0)=1>0,
          \qquad
          f(1)=1-3+1=-1<0 .
        \]
        The sign changes, so by the IVT there is an $x^{*}\in(0,1)$ with
        $f(x^{*})=0$. \textbf{Bisection step:} midpoint $m=\tfrac12$,
        \[
          f\Bigl(\tfrac12\Bigr)
          =\tfrac18-\tfrac32+1=-\tfrac38<0 .
        \]
        The sign change now sits between $f(0)>0$ and
        $f(\tfrac12)<0$: the new bracket is $\bigl[0,\tfrac12\bigr]$,
        half as long as before.

  \item The numerator $\sqrt{x+2}$ is the composition of the continuous
        functions $x\mapsto x+2$ and $t\mapsto\sqrt{t}$, hence
        continuous wherever it is defined, i.e.\ for $x+2\ge 0$. The
        denominator $x-1$ is continuous everywhere, and a quotient of
        continuous functions is continuous wherever the denominator is
        nonzero. Therefore $g$ is continuous on its whole domain
        \[
          [-2,1)\cup(1,\infty),
        \]
        i.e.\ for all $x\ge -2$ with $x\neq 1$.
\end{solutions}

\section{Additional Problems}

\begin{problems}
  \item $\displaystyle\lim_{x\to -2}\frac{x^{2}-x-6}{x+2}$
  \item $\displaystyle\lim_{x\to 9}\frac{\sqrt{x}-3}{x-9}$
  \item $\displaystyle\lim_{x\to 0}\frac{\sqrt{4+x}-2}{x}$
  \item $\displaystyle\lim_{x\to\infty}\frac{5x^{3}+2x}{2x^{3}-x^{2}+4}$
        \quad and \quad
        $\displaystyle\lim_{x\to-\infty}\frac{5x^{3}+2x}{2x^{3}-x^{2}+4}$
  \item Find all asymptotes of $f(x)=\dfrac{3x+1}{x-2}$, including the
        one-sided behaviour at the vertical asymptote.
  \item Find the constant $c$ that makes
        \[
          f(x)=
          \begin{cases}
            cx^{2}, & x\le 2,\\
            3x+2,   & x>2,
          \end{cases}
        \]
        continuous on $\R$.
  \item Classify all discontinuities of
        $f(x)=\dfrac{x^{2}-1}{x^{2}-3x+2}$.
  \item Use the sandwich rule to compute
        $\displaystyle\lim_{x\to\infty}\frac{2+\cos x}{x}$.
  \item Show that $f(x)=x^{3}+x-3$ has a root in $[1,2]$ and perform
        one bisection step, stating the new bracket.
  \item Explain why $f(x)=x^{2}-4x+1$ attains a maximum and a minimum
        on $[0,3]$, and find both values (no derivatives: complete the
        square).
\end{problems}

\section{Additional Solutions}

\begin{solutions}
  \item Form $0/0$. Factor $x^{2}-x-6=(x-3)(x+2)$:
        \[
          \lim_{x\to -2}\frac{(x-3)(x+2)}{x+2}
          =\lim_{x\to -2}\,(x-3)=-5 .
        \]

  \item Form $0/0$. Conjugate $\sqrt{x}+3$, using
        $x-9=(\sqrt{x}-3)(\sqrt{x}+3)$:
        \[
          \frac{\sqrt{x}-3}{x-9}=\frac{1}{\sqrt{x}+3}
          \;\xrightarrow[x\to 9]{}\;
          \frac{1}{3+3}=\frac{1}{6} .
        \]

  \item Form $0/0$. Conjugate $\sqrt{4+x}+2$:
        \[
          \frac{\sqrt{4+x}-2}{x}
          =\frac{(4+x)-4}{x\bigl(\sqrt{4+x}+2\bigr)}
          =\frac{1}{\sqrt{4+x}+2}
          \;\xrightarrow[x\to 0]{}\;
          \frac{1}{4} .
        \]

  \item Equal degrees: the limit is the ratio of the leading
        coefficients,
        \[
          \lim_{x\to\pm\infty}\frac{5x^{3}+2x}{2x^{3}-x^{2}+4}
          =\frac{5}{2}
        \]
        in both directions (divide numerator and denominator by
        $x^{3}$).

  \item \textbf{Horizontal:} equal degrees give
        $y=\tfrac{3}{1}=3$ as $x\to\pm\infty$. \textbf{Vertical:} the
        denominator vanishes at $x=2$, the numerator equals $7\neq 0$
        there, so $x=2$ is a pole. Near $x=2$ the numerator is close to
        $7>0$, so the sign follows $x-2$:
        \[
          \lim_{x\to 2^{-}}f(x)=-\infty,
          \qquad
          \lim_{x\to 2^{+}}f(x)=+\infty .
        \]
        Asymptotes: $y=3$ and $x=2$.

  \item Both branches are continuous; match them at the seam $x=2$:
        \[
          \lim_{x\to 2^{-}}cx^{2}=4c=f(2),
          \qquad
          \lim_{x\to 2^{+}}(3x+2)=8 .
        \]
        Continuity requires $4c=8$, so $c=2$.

  \item Factor: $x^{2}-1=(x-1)(x+1)$ and $x^{2}-3x+2=(x-1)(x-2)$, so
        for $x\neq 1$
        \[
          f(x)=\frac{(x-1)(x+1)}{(x-1)(x-2)}=\frac{x+1}{x-2}.
        \]
        \textbf{At $x=1$:} the factor cancels,
        $\lim_{x\to 1}f(x)=\frac{2}{-1}=-2$ exists while $f(1)$ is
        undefined: \textbf{removable}. \textbf{At $x=2$:} the numerator
        tends to $3\neq 0$ while the denominator vanishes:
        \[
          \lim_{x\to 2^{-}}\frac{x+1}{x-2}=-\infty,
          \qquad
          \lim_{x\to 2^{+}}\frac{x+1}{x-2}=+\infty,
        \]
        a \textbf{pole} with vertical asymptote $x=2$.

  \item Since $-1\le\cos x\le 1$, we have $1\le 2+\cos x\le 3$, and for
        $x>0$
        \[
          \frac{1}{x}\;\le\;\frac{2+\cos x}{x}\;\le\;\frac{3}{x} .
        \]
        Both bounds tend to $0$ as $x\to\infty$, so the limit is $0$ by
        the sandwich rule.

  \item $f(x)=x^{3}+x-3$ is continuous (polynomial), and
        \[
          f(1)=1+1-3=-1<0,
          \qquad
          f(2)=8+2-3=7>0,
        \]
        so by the IVT a root lies in $(1,2)$. \textbf{Bisection step:}
        midpoint $m=1.5$,
        \[
          f(1.5)=3.375+1.5-3=1.875>0 .
        \]
        The sign change is now between $f(1)<0$ and $f(1.5)>0$: new
        bracket $[1,\,1.5]$.

  \item $f$ is a polynomial, hence continuous, and $[0,3]$ is a closed
        bounded interval, so by the Weierstrass theorem $f$ attains a
        maximum and a minimum on $[0,3]$. Completing the square,
        \[
          f(x)=x^{2}-4x+1=(x-2)^{2}-3 ,
        \]
        which is smallest when $(x-2)^{2}=0$: minimum $-3$ at $x=2$
        (which lies in $[0,3]$). The maximum of $(x-2)^{2}$ over
        $[0,3]$ occurs at the endpoint farther from $2$, namely $x=0$:
        $f(0)=1$ versus $f(3)=9-12+1=-2$. Maximum $1$ at $x=0$,
        minimum $-3$ at $x=2$.
\end{solutions}


\end{document}
